% v4.10.0 post-release layers (source-only research layers after the v4.10.0 release, 2026-11-06).
% Layers are appended here; earlier layers are never rewritten.
% ============================================================================
\clearpage
\section*{v4.10.0-r01 Source-Only: The Full Einstein Equation of W1 and the Sign of the Gravitational Action}
\addcontentsline{toc}{section}{v4.10.0-r01 Source-Only: The Full Einstein Equation of W1 and the Sign of the Gravitational Action}
\label{sec:v4100-r01-einstein-sign}
% ============================================================================

\paragraph{Purpose.}
Layer v4.9.0-r86 derived the field equations of the joint action W1 (\S\ref{sec:v490-r81-joint-action-w1}) from their definitions up to
the full non-linear Einstein equation and the Bianchi identity (\S\ref{sec:v490-r86-wolfram-checks}).  This first layer after the
v4.10.0 release closes that gap with one pre-registered test (H90, not blind) and records what it found on the way: W1 as written has the
wrong sign of the gravitational action.  The author decided on 6 November 2026 to correct it.  The layer states the corrected action and
lists which earlier statements change sign.  Earlier layers are not rewritten.  The notes are
\texttt{notes/\allowbreak research/\allowbreak v4\_10\_0\_r01/}.

\noindent\textbf{Source-only status.}
This v4.10.0-r01 layer is source-only.  It does not build a PDF, does not assign a new DOI,
and does not authorize a public release.  The v4.10.0 Version DOI
\href{https://doi.org/10.5281/zenodo.23991357}{10.5281/zenodo.23991357} and the Concept DOI
\href{https://doi.org/10.5281/zenodo.19210728}{10.5281/zenodo.19210728} are unchanged by it.

\subsection*{v4.10.0-r01 results}
\label{subsec:v4100-r01-results}

\paragraph{Theoremlet v4.10.0-r01.1 (H90: the Einstein--Cartan equations of W1, non-linear).}
Pre-registered (commit \texttt{e5d03d8}, not blind).  Conventions as in r86: \(\eta=\operatorname{diag}(1,-1,-1,-1)\),
\(D_\mu\psi=\partial_\mu\psi+\tfrac12\omega_{\mu ab}\Sigma^{ab}\psi+iqJA_\mu\psi\),
\(R^{ab}{}_{\mu\nu}=\partial_\mu\omega_\nu{}^{ab}-\partial_\nu\omega_\mu{}^{ab}+\omega_\mu{}^a{}_c\,\omega_\nu{}^{cb}-\omega_\nu{}^a{}_c\,\omega_\mu{}^{cb}\), and the gravitational action
written as \(S_{\rm grav}=\tfrac{\sigma_g}{2\kappa}\int e\,R\) (W1 as written in r81.1: \(\sigma_g=+1\)).  At a generic point of three random curved
backgrounds (vierbein a random polynomial of degree three, not diagonal) all fields were carried as jets to third order: three Weyl
fields with \(J=(+1,-1,+1)\), a real potential \(A\), the vierbein.  The spinor jets solve the Weyl equation with gauge coupling and
four-fermion term at the point and to first derivative, \(A\) the Maxwell equation at the point.  Both signs \(\sigma_g=\pm1\) were computed.
\begin{enumerate}[leftmargin=2em]
\item[(a)] Cartan equation.  The spin connection \(\tilde\omega\) of the tetrad postulate solves the full \(\omega\)-equation (with its derivative
term) without spinors; with spinors \(\omega=\tilde\omega+K\) solves it, with \(K_{cab}=-\tfrac{\sigma_g}{4}\kappa\,\varepsilon_{cabd}\,j^d\) totally
antisymmetric (residuals \(\le2\times10^{-14}\)).
\item[(b)] Einstein--Cartan equation.  The variation with respect to the vierbein at fixed \(\omega=\tilde\omega+K\) is symmetric on the field
equations and equals a multiple of \(\tilde G_{\mu\nu}+\sigma_g\kappa\,T^{\rm eff}_{\mu\nu}\) (residual \(10^{-15}\)), with
\[
T^{\rm eff}_{\mu\nu}=T^{\rm Bel}_{\mu\nu}-g_{\mu\nu}\mathcal L_W+T^{\rm Maxwell}_{\mu\nu}-g_{\mu\nu}\,\lambda\,j\cdot j,
\qquad \lambda=-\sigma_g\,\tfrac{3\kappa}{16}.
\]
Hence \(\tilde G=-\sigma_g\,\kappa\,T^{\rm eff}\), fully non-linear.  The \(K\) term of the Einstein--Cartan Weyl equation equals
\(2\lambda\,j_a\bar\sigma^a\psi\).
\item[(c)] Compatibility.  \(\tilde\nabla^\mu T^{\rm eff}_{\mu\nu}=0\) on the field equations (relative \(4\times10^{-15}\)); \(\tilde\nabla^\mu\tilde G_{\mu\nu}=0\) as a
control.  Negative controls: without the four-fermion term the residuals rise to \(0.12\)--\(0.19\) (Einstein) and \(0.25\)--\(0.81\)
(divergence); with spinors off the field equations to \(0.16\)--\(0.19\) and \(0.92\)--\(1.05\).
\item[(d)] Sign.  For W1 as written (\(\sigma_g=+1\)) the factor is \(-1\) to \(4\times10^{-15}\): \(\tilde G=-\kappa T^{\rm eff}\).  A second system
(Wolfram) gives \(G_{00}=2\nabla^2\Phi\) for \(g=\operatorname{diag}(1+2\Phi,-(1-2\Phi)\mathbb 1)\) in the same curvature convention, so
\(\tilde G=+\kappa T\) is attraction and \(\tilde G=-\kappa T\) repulsion.  W1 as written repels positive energy.
\end{enumerate}
The Hehl--Datta control of r81.1(b) compared only the modulus \(3/16\); the sign had not been checked.

\paragraph{Decision v4.10.0-r01.2 (author, 6 November 2026: the corrected action).}
W1 is corrected to
\[
S=\int e\,\Bigl[-\tfrac1{2\kappa}R(e,\omega)+\tfrac12\bigl(\Psi^\dagger i\bar\sigma^\mu D_\mu\Psi+\mathrm{c.c.}\bigr)-\tfrac14F^2\Bigr]
\]
in the signature \((+,-,-,-)\) (as in Landau--Lifshitz).  The field equations become \(\tilde G=\kappa\,T^{\rm eff}\) with
\(\mathcal L_{4f}=+\tfrac3{16}\kappa\,j\cdot j\), \(T^{(4f)}_{\mu\nu}=-\tfrac{3\kappa}{16}\,j^2g_{\mu\nu}\), Weyl term \(+\tfrac{3\kappa}{8}j_a\bar\sigma^a\psi\),
\(K_{cab}=+\tfrac\kappa4\varepsilon_{cabd}j^d\).

\paragraph{Correction v4.10.0-r01.3 (statements that change sign).}
Moduli and forms stay; the signs change as follows.
\begin{enumerate}[leftmargin=2em]
\item[(a)] r81.1(b) and r86.2(b), (c): \(\mathcal L_{\rm eff}=+\tfrac3{16}\kappa\,j\cdot j\) and \(+\tfrac{3\kappa}{8}j_a\bar\sigma^a\psi\) in the Weyl equation (were
\(-\)); correct for the action as written, not for the physical one.
\item[(b)] r86.3(a), (b): \(T^{(4f)}=-\tfrac{3\kappa}{16}j^2g\) (\(\Lambda\) form, negative); the homogeneous state of two opposite spins has
\(\rho=p=-\tfrac{3\kappa}{16}n^2<0\), the negative spin energy density known from Einstein--Cartan theory.
\item[(c)] Correction r86.4 is withdrawn in its sign part: the original reading of the W1 notes, \(\rho\propto-\kappa n^2\), holds for the physical
sign.  ``Not like \(\Lambda\)'' for the whole tensor and r81.1(c) (no vacuum term) stand.
\item[(d)] r82.4: the pairing channel of the torsion interaction is repulsive, \(H_{\rm int}\supset+\tfrac34\kappa\,\Phi^\dagger\Phi\).  No condensate arises
through it; the \(\nu^c\) Majorana mass near \(M_P\) and the seesaw estimate of r82.4 fall.
\item[(e)] Not affected (review in \texttt{llm/\allowbreak Durchsicht\_\allowbreak Vorzeichen\_\allowbreak kappa\_\allowbreak r81\_\allowbreak r86.md}): H69--H73 (source side only),
H74 (no gravity), H75 (computed with Newtonian attraction), H76 (rate \(\propto\kappa^4\)), H77--H81 and H84--H86 (no gravitational coupling), H82
and H83 (relative weights and phases of the Einstein--Cartan term), H48 (Hehl--Datta sign taken from the literature).
\end{enumerate}

\paragraph{Proof and measurement sources.}
Script \texttt{k205\_\allowbreak w1\_\allowbreak einstein\_\allowbreak nichtlinear.py} (Python, jets, output \texttt{k205\_\allowbreak out.json}) and Wolfram code
\texttt{k205b\_\allowbreak schwachfeld.wl}; notes in \texttt{notes/\allowbreak research/\allowbreak v4\_10\_0\_r01/}; pre-registration and result H90 in
\texttt{notes/\allowbreak research/\allowbreak predictions/}.  \(\square\)

\paragraph{Consequence.}
The field equations of W1 are now derived independently in full: Cartan equation, elimination of the torsion, Einstein--Cartan equation, its
Riemannian form with the four-fermion term, and its compatibility with the Bianchi identity.  With the corrected sign the four-fermion
part has the form of a negative vacuum term with a height \(\propto j^2\); the whole tensor of a homogeneous state is stiff matter with
negative energy, not \(\Lambda\).  W1 still produces no vacuum term.  OP5 is not closed.

\subsection*{v4.10.0-r01 status in the three frames}
\label{subsec:v4100-r01-frames}

\begin{itemize}[leftmargin=2em]
\item r01.1--r01.3 concern the matter part of W1 in frame M.  The tachyon image under the second reading (r84) multiplies the action by the
weights \(a\), \(b\); the sign correction is common to both and does not change their ratio.  The photon frame P is the degenerate fixed point;
no frame-free statement is made beyond that.
\item r01.3(d) concerns the pairing of two left-handed fields in frame M; under the second reading the image term is \(-\mathcal L_M\) (r84), so the
interaction keeps its sign relative to the kinetic term.
\end{itemize}

\subsection*{v4.10.0-r01 register entries}
\label{subsec:v4100-r01-register}

\begin{center}
\small
\begin{tabular}{@{}>{\raggedright\arraybackslash}p{0.34\textwidth}>{\raggedright\arraybackslash}p{0.58\textwidth}@{}}
\toprule
statement & status after v4.10.0-r01\\
\midrule
full Einstein equation of W1, Bianchi & derived and compatible (H90)\\
sign of the gravitational action & corrected to \(-\tfrac1{2\kappa}\int eR\) (author decision)\\
four-fermion term & \(+\tfrac3{16}\kappa\,j\cdot j\); \(T^{(4f)}\) negative \(\Lambda\) form\\
homogeneous spin state & \(\rho=p=-\tfrac{3\kappa}{16}n^2<0\)\\
torsion pairing channel (r82.4) & repulsive; Planck-scale \(\nu^c\) condensate withdrawn\\
W1 vacuum term & still absent\\
\bottomrule
\end{tabular}
\end{center}

\subsection*{v4.10.0-r01 ledger and guard update}
\label{subsec:v4100-r01-ledger-guard}

\[
R_{\rm chk}^{v4100\text{-}r01}=(1,1,1)
\]
for (Einstein--Cartan equations derived; compatibility shown; sign corrected by author decision).

The layer spends no guard:
\[
R_{\rm glob}^{(v4100\text{-}r01,+)}=(0,1,1,1,1),
\qquad
\Gamma_{\rm guard}^{v4100\text{-}r01}=(1,1,1,1,1),
\qquad
\beta_4=1/256\ \text{unspent}.
\]
It does not imply exact finite \(C5\) zero closure, exact all-readout terminality,
\(C6/T2\), a global photon-clock/RH theorem, or global all-readout closure.

% ============================================================================
\clearpage
\section*{v4.10.0-r02 Source-Only: Asymptotics of the Admixture at the Big Bang}
\addcontentsline{toc}{section}{v4.10.0-r02 Source-Only: Asymptotics of the Admixture at the Big Bang}
\label{sec:v4100-r02-admixture-asymptotics}
% ============================================================================

\paragraph{Purpose.}
Under the domain choice G2 the big bang is a transition from the tachyon reading to the matter reading, and a mode of the tachyon
reading picks up an admixture \(C\) of positive frequency after the bang that falls off for short waves (r85.6, \S\ref{sec:v490-r85-masses-cpt}).
Layer v4.9.0-r86 found the fall-off close to \(e^{-\pi k\Delta}\) with an unexplained power prefactor (\S\ref{sec:v490-r86-wolfram-checks}).
This layer records three pre-registered tests on that question (H91--H93).  The last one derives the asymptotics in closed form.  All three have
formal defects in their validity conditions, which are disclosed.  The notes are
\texttt{notes/\allowbreak research/\allowbreak v4\_10\_0\_r02/}.  Earlier layers are not rewritten.

\noindent\textbf{Source-only status.}
This v4.10.0-r02 layer is source-only.  It does not build a PDF, does not assign a new DOI,
and does not authorize a public release.  The v4.10.0 Version DOI
\href{https://doi.org/10.5281/zenodo.23991357}{10.5281/zenodo.23991357} and the Concept DOI
\href{https://doi.org/10.5281/zenodo.19210728}{10.5281/zenodo.19210728} are unchanged by it.

\subsection*{v4.10.0-r02 results}
\label{subsec:v4100-r02-results}

Throughout: \((w\chi')'+k^2w\chi=0\), \(w=e^{Ls}\), \(s=\tfrac12(1+\tanh(t/\Delta))\), \(L=\ln(a/b)\mp i\pi\) for the two senses of rotation, \(a,b\) the
matter weights of r84, \(\Delta=2\); \(C\) is the coefficient of \(e^{-ikt}\) after the bang for \(e^{ikt}\) incoming before it.

\paragraph{Theoremlet v4.10.0-r02.1 (H91: the fall-off up to \(k\Delta=60\)).}
Pre-registered (commit \texttt{a49ed7e}, partly not blind).  Integration along \(t=x+ih\) below the singularity; two contours agree to
\(10^{-10}\) up to \(k\Delta=60\).
\begin{enumerate}[leftmargin=2em]
\item[(a)] Invalid by its own criterion, which was a planning error: agreement with the H87 values to \(10^{-5}\) was required, but H87 had recorded
four digits only.  All deviations, up to \(3.2\times10^{-4}\), lie within that rounding.  By the pre-registration parts (b)--(d) are then invalid.  As
formulated they would also have missed their criteria.
\item[(b)--(d)] A fit with \(\sqrt{k\Delta}\) and \(\ln k\Delta\) gives a slope \(3.049\), \(3\%\) below \(\pi\).  The \(\sqrt{k\Delta}\) term improves the fit only by a
factor \(6.5\).  The slope between \(50\) and \(60\) is \(3.021\).
\end{enumerate}
\emph{Noticed afterwards.}  The local slope is wavy and rises slowly (\(3.026\) at \(60\)).  \(|C|\) is the same for both senses of rotation, to
\(10^{-10}\); the asymmetry of r85.6(d) concerns only the other coefficient.

\paragraph{Theoremlet v4.10.0-r02.2 (H92: the prefactor is the Stokes function of the local problem).}
Pre-registered (commit \texttt{c347a92}, partly not blind).  Near \(t_0=i\pi\Delta/2\), \(s\approx\tfrac12+\Delta/(2\tau)\) with \(\tau=t-t_0\), so \(w\) has an
essential singularity.  In \(\sigma=k\tau\) the local equation is \(y''-(\lambda/\sigma^2)y'+y=0\), with the single parameter \(\lambda=Lk\Delta/2\).  Its
reflection \(r(\lambda)\), continued below \(\sigma=0\), gives \(C\approx r(\lambda)e^{2ikt_0}=r(\lambda)e^{-\pi k\Delta}\).
\begin{enumerate}[leftmargin=2em]
\item[(b)] Confirmed: \(\ln|C|+\pi k\Delta-\ln|r|\) is smooth; the fit residual on \([20,60]\) is \(4.1\times10^{-3}\).  The waves sit in \(r\).
\item[(c)] Confirmed: \(|r(\bar\lambda)|=|r(\lambda)|\) to \(8\times10^{-11}\).
\item[(d)] Confirmed: \(r(\lambda)/(2\pi\lambda)=1-1.6\times10^{-8}\) at \(|\lambda|=10^{-3}\), the Born limit.
\item[(a)] Missed: \(\ln|C|+\pi k\Delta-\ln|r|\to\approx-0.62\), not \(0\).
\end{enumerate}
\emph{Disclosure.}  In run 1 the accuracy control (a contour at \(1.5h\)) lost digits (\(3.7\)--\(5.5\times10^{-6}\), above \(10^{-6}\)).  In run 2 the control
used \(0.7h\) (\(\le6\times10^{-10}\)); the values of \(r\) are identical in both runs.

\paragraph{Theoremlet v4.10.0-r02.3 (H93: the asymptotics derived).}
Pre-registered with the derivation (commit \texttt{a4b1d17}, partly not blind).
\begin{enumerate}[leftmargin=2em]
\item[(i)] With \(\chi=w^{-1/2}u\), \(u''+Qu=0\) and \(w\) going from \(1\) to \(e^L\): \(C=e^{-L/2}C_u\).  This is the factor missing in r02.2(a).
\item[(ii)] For large \(\lambda\) the local \(Q=1-\lambda^2/(4\sigma^4)-\lambda/\sigma^3\) has four turning points.  The over-barrier formula
\(r\approx-i\,e^{2iF}\) at the turning point \(\eta=-i\) (in \(\sigma=\sqrt{\lambda/2}\,\eta\)) gives \(\ln|r|=c\sqrt{k\Delta}\) with
\[
c=\frac{2\sqrt\pi\,\Gamma(3/4)}{\Gamma(1/4)}\,\operatorname{Re}\sqrt L=1.809205,
\]
\[
1+\int_1^\infty\bigl(1-\sqrt{1-\eta^{-4}}\bigr)d\eta=\frac{2\sqrt\pi\,\Gamma(3/4)}{\Gamma(1/4)}=1.19814.
\]
The term \(-\lambda/\sigma^3\) contributes only the phase \(-\pi/4\), so no additive constant appears.
\item[(iii)] Together:
\[
\ln|C|=-\pi k\Delta+c\sqrt{k\Delta}-\tfrac12\ln(a/b)+O\bigl((k\Delta)^{-1/2}\bigr).
\]
\end{enumerate}
Numbers (run 2): slope of \(\ln|r|\) between \(k\Delta=300\) and \(600\) is \(1.808965\) (\(-0.013\%\), criterion \(0.3\%\)); \(\ln|r|-c\sqrt{k\Delta}=0.0057,\ 0.0040\)
(criterion \(\le0.05\), falling); the full problem at \(k\Delta=80,100,120\) deviates from (iii) by \(-0.0028,\ -0.0085,\ -0.0045\) (criterion \(0.02\)).  Both
senses of rotation give the same numbers.
\emph{Disclosure.}  The validity condition (accuracy controls \(\le10^{-6}\)) is not met in one place.  For the sense \(+\) at \(k\Delta=600\) the control is
\(1.4\times10^{-5}\) in run 2 (contour \(1.3h\) instead of \(0.7h\), changed after run 1, where the \(0.7h\) contour lost digits up to \(1.0\)).  Formally the
test is therefore invalid.  For the sense \(-\) it is valid in full and confirmed, and every content criterion is met with a margin of a factor
\(2\) to \(20\).

\paragraph{Proof and measurement sources.}
Scripts \texttt{k206\_\allowbreak abfall\_\allowbreak wurzelterm.py}, \texttt{k207\_\allowbreak lokales\_\allowbreak problem.py},
\texttt{k208\_\allowbreak wkb\_\allowbreak pruefung.py} with their outputs, and the notes in \texttt{notes/\allowbreak research/\allowbreak v4\_10\_0\_r02/};
pre-registrations and results H91--H93 in \texttt{notes/\allowbreak research/\allowbreak predictions/}.  \(\square\)

\paragraph{Consequence.}
The admixture at the big bang now has a closed asymptotic form.  The exponent \(\pi\) comes from the position of the singularity of \(\tanh\), the
\(\sqrt{k\Delta}\) term from the essential singularity of the weight \(e^{Ls}\) there, and the constant from the adiabatic amplitude between the
weights before and after the bang.  This explains the slope below \(\pi\) of r86.1(b), the waves (from the subdominant turning point, not checked
separately) and the missing factor of r02.2(a).  The cut-off scale stays about \(1/(\pi\Delta)\), shifted only by a factor growing like
\(e^{1.8\sqrt{k\Delta}}\); \(\Delta\) remains free.  OP5 is not closed.

\subsection*{v4.10.0-r02 status in the three frames}
\label{subsec:v4100-r02-frames}

\begin{itemize}[leftmargin=2em]
\item The transition joins frame T (before the bang) and frame M (after it) under G2, with the weights \(a\), \(b\) of frame M.  All results concern
this M--T transition.  A frame-free statement is not made.
\item Photon frame P: with photon weights \(a=b\) one has \(L=\mp i\pi\), the constant \(-\tfrac12\ln(a/b)\) vanishes and \(c\) would be
\(1.19814\sqrt{\pi/2}=1.502\).  This is a consequence noticed afterwards and was not computed.  In the sharp limit the photon weights give no
admixture (r85.5), unchanged.
\end{itemize}

\subsection*{v4.10.0-r02 register entries}
\label{subsec:v4100-r02-register}

\begin{center}
\small
\begin{tabular}{@{}>{\raggedright\arraybackslash}p{0.34\textwidth}>{\raggedright\arraybackslash}p{0.58\textwidth}@{}}
\toprule
statement & status after v4.10.0-r02\\
\midrule
fall-off up to \(k\Delta=60\) & numbers secured; H91 invalid (planning error)\\
prefactor as local Stokes function & waves, symmetry and Born limit confirmed; normalization missed (H92)\\
closed asymptotics of \(\ln|C|\) & derived; numbers met; formally invalid for sense \(+\), valid and confirmed for sense \(-\) (H93)\\
\(|C|\) for both senses & equal (afterwards, H91; local form H92c)\\
cut-off scale & \(\approx1/(\pi\Delta)\), \(\Delta\) free\\
\bottomrule
\end{tabular}
\end{center}

\subsection*{v4.10.0-r02 ledger and guard update}
\label{subsec:v4100-r02-ledger-guard}

\[
R_{\rm chk}^{v4100\text{-}r02}=(1,1,0)
\]
for (fall-off computed to \(k\Delta=600\); asymptotics derived; one accuracy control above its bound).

The layer spends no guard:
\[
R_{\rm glob}^{(v4100\text{-}r02,+)}=(0,1,1,1,1),
\qquad
\Gamma_{\rm guard}^{v4100\text{-}r02}=(1,1,1,1,1),
\qquad
\beta_4=1/256\ \text{unspent}.
\]
It does not imply exact finite \(C5\) zero closure, exact all-readout terminality,
\(C6/T2\), a global photon-clock/RH theorem, or global all-readout closure.

% ============================================================================
\clearpage
\section*{v4.10.0-r03 Source-Only: Fermions and Bosons at the Big Bang, and the Vacuum Term under G2}
\addcontentsline{toc}{section}{v4.10.0-r03 Source-Only: Fermions and Bosons at the Big Bang, and the Vacuum Term under G2}
\label{sec:v4100-r03-bang-g2}
% ============================================================================

\paragraph{Purpose.}
Layer v4.10.0-r02 closed the asymptotics of the bosonic admixture at the big bang (\S\ref{sec:v4100-r02-admixture-asymptotics}).
This layer records what follows for the other fields and for the vacuum term under the domain choice G2 (tachyon reading before,
matter reading after the bang; r85).  It contains one pre-registered test (H94), an author decision on the duration \(\Delta\) of the
transition, an exact symmetry of the two senses of rotation, and analytic explorations of four routes to \(\Lambda\) under G2.  Two of my own
intermediate claims are corrected here.  The notes are
\texttt{notes/\allowbreak research/\allowbreak v4\_10\_0\_r03/}.  Earlier layers are not rewritten.

\noindent\textbf{Source-only status.}
This v4.10.0-r03 layer is source-only.  It does not build a PDF, does not assign a new DOI,
and does not authorize a public release.  The v4.10.0 Version DOI
\href{https://doi.org/10.5281/zenodo.23991357}{10.5281/zenodo.23991357} and the Concept DOI
\href{https://doi.org/10.5281/zenodo.19210728}{10.5281/zenodo.19210728} are unchanged by it.

\subsection*{v4.10.0-r03 results}
\label{subsec:v4100-r03-results}

\paragraph{Theoremlet v4.10.0-r03.1 (fermions pass the bang without admixture).}
Exploration (analytic, no criterion).  Under G2 and the second reading the symmetric Dirac or Weyl action is weighted by \(c(\tau)\), going from
\(-b\) to \(a\).  The field equation \(c(i\gamma^\mu\partial_\mu-m)\psi+\tfrac i2c'\gamma^0\psi=0\) is solved exactly by \(\psi=c^{-1/2}\chi\), with \(\chi\) free, for
every profile (sympy, residual \(0\)).
\begin{enumerate}[leftmargin=2em]
\item[(a)] Positive and negative frequencies do not mix: no fermions are created at the bang, for any \(k\) and any profile.  A first-order equation lets
the weight be scaled away; the second-order equations of scalars and gauge fields do not.
\item[(b)] Transfer factor \(\psi_{\rm after}/\psi_{\rm before}=\pm i\sqrt{b/a}=\pm i\beta_{\rm sym}\) (continuous branch along the path; the two senses differ in phase
only).  The map is unitary in the weighted norm, and the frequency is kept: a particle of the T reading becomes an antiparticle-like state of the M
reading, exactly and for all \(k\).
\item[(c)] Correction (made in the notes the same day): mass terms keep their sign under the image (r85, H82), so before the bang the mass is \(-m\)
relative to the kinetic term.  On FRW the mass term is \(m\,a(\tau)\) with \(a\to0\) at the bang, so (a) and (b) hold for \(k\gg m\,a\) at the transition.
\end{enumerate}

\paragraph{Theoremlet v4.10.0-r03.2 (H94: the weighted norm at the bang).}
Pre-registered (commit \texttt{99f1cce}, partly not blind).  \(N=a(|C|^2-|D|^2)/b\) for the mode of positive norm before the bang (tanh profile,
\(\Delta=2\)).  All four parts confirmed (one run, Wronskian constant to \(10^{-11}\)).
\begin{enumerate}[leftmargin=2em]
\item[(a)] Bosons: \(N\to+1\) for long waves (\(1.024\), \(0.974\) at \(k\Delta=0.01\)) and \(N\to-1\) for short waves (\(-0.969\), \(-1.032\) at \(20\)).
\item[(b)] Exactly one sign change per sense of rotation, at \(k^*\Delta=0.7361\) and \(0.2569\).
\item[(c)] In between the norm is not conserved (\(0.1<|N|<0.9\) at \(10\) and \(5\) grid points).
\item[(d)] Fermions: \(N_F=-1\) on all points, to \(6\times10^{-12}\).
\end{enumerate}
\emph{Noticed afterwards.}  The bosonic norm is also amplified (up to \(1.18\) and \(-1.63\)).  The mean of the two senses does not restore it.  The
statement ``bosons with matter weights are created, \(n=0.404\) per mode'' holds only in the picture of the sharp real jump.

\paragraph{Decision v4.10.0-r03.3 (author, 6 November 2026: \(\Delta\lesssim\ell_P\)).}
The duration of the transition is taken at or below the Planck length.  All modes below the Planck scale then have \(k\Delta\ll1\): bosons with matter
weights pass norm-preserving as a Bogoliubov transformation with \(n=(1-\beta_{\rm sym}^2)^2/(4\beta_{\rm sym}^2)=0.404\) per mode, and fermions are mirrored
by CPT without creation.  The non-unitary range \(k\sim1/\Delta\) lies at or beyond the Planck scale.  The total energy of the created bosons is set by
that range and stays open.

\paragraph{Theoremlet v4.10.0-r03.4 (the two senses of rotation are exactly conjugate).}
Exploration (derived analytically, checked numerically).  For the transfer matrix \(M(L)=\begin{pmatrix}a&b\\c&d\end{pmatrix}\) in the basis
\((e^{ikt},e^{-ikt})\) one has \(\det M=e^{-L}\) (weighted Wronskian).  The duality \(v=wy'\) gives \(M(-L)=e^L JMJ\) with \(J=\operatorname{diag}(1,-1)\).  Reflection
with conjugation, using \(s(-t)=1-s(t)\), gives \(M(-\bar L)=\overline{M(L)}^{-1}\).  Together
\[
M(\bar L)=\begin{pmatrix}\bar d&\bar b\\ \bar c&\bar a\end{pmatrix}.
\]
Hence \(C(\bar L)=\overline{C(L)}\), so \(|C_+|=|C_-|\) exactly, and \(D_+\overline{D_-}-b_+C_+=e^{-L}\), the product rule of r85.6 with its correction term.  For
the local problem of r02.2, \(M=\begin{pmatrix}1&0\\r&1\end{pmatrix}\) exactly, so \(r(\bar\lambda)=\overline{r(\lambda)}\).  Numerically all identities hold to
\(\le3\times10^{-10}\) at \(k\Delta=0.01\)--\(5\).

\paragraph{Theoremlet v4.10.0-r03.5 (no \(\Lambda\) from the tachyon sector after the bang under G2).}
Exploration (analytic, no criterion).  Four routes were examined.
\begin{enumerate}[leftmargin=2em]
\item[(a)] Mixture (deferred idea I3): its \(\Lambda\) side came from the tension of the T reading.  Under G2 only the M reading exists after the bang, so
only the dust part remains.
\item[(b)] Tachyonic condensation (I4): after the bang there are no tachyon modes.  Before it the image of a scalar sees \(-V\) (\(\mathcal L_T=-\tfrac12X-V\),
r84); a potential stable in M runs away there (\(\ddot\phi=m^2\phi+\lambda\phi^3\)) instead of condensing.  Under G1 one has \(m_{\rm eff}^2=\gamma m^2>0\).
\item[(c)] Ghost condensate (r83.4): under G1, \(aP_M+bP_T=(a-b)X/2+(a+b)cX^2\) is healthy and does not condense.  Under G2 the condensate exists only
before the bang; after it \(P_M\) has \(w\in[\tfrac13,1]\), never \(-1\); at the bang the condensate has \(\pi=0\), so the M field is at rest.
\item[(d)] Conclusion: under G2 a vacuum term after the bang can only be \(aV_{\min}\) of the M reading, with a free value.  The structure does not fix
\(\Lambda\) under G2.
\end{enumerate}

\paragraph{Theoremlet v4.10.0-r03.6 (boundary condition and spatial reflection; correction).}
Exploration (analytic).  If the T field before the bang is identified pointwise with the image of the M field after it,
\(\phi_T(x)=\phi_M(\Lambda(i\pi)x)\), the junction yields \(\phi'(0)=0\) for scalars and \(\psi(0)=0\) for fermions (\(\beta_{\rm sym}\neq1\)).  The same identification
needs a direction \(n\) and a spatial origin (from \(x\mapsto\Lambda x+2x_0\)) that nothing fixes; classically it would correlate the modes \(k\) and \(P_nk\).
Quantum mechanically it is inconsistent for \(a\neq b\): with antiunitary time reversal the junction forces \(u'(0)=0\), and the Wronskian
normalization \(uu'^*-u^*u'=i\) then fails at \(\tau=0\).  \emph{Correction:} the conditions \(\phi'(0)=0\), \(\psi(0)=0\) stated the same day in the notes do
not follow from the structure; a frozen scalar at the bang is a possible choice of state, not a consequence.  Under G2 the second reading is a
statement about the action, \(-b\mathcal L_M\) before and \(a\mathcal L_M\) after, with independent data before the bang, as in H85--H94; at that level
neither a direction nor an origin appears (r84, H79).

\paragraph{Proof and measurement sources.}
Scripts \texttt{k209\_\allowbreak fermionen\_\allowbreak urknall.py}, \texttt{k210\_\allowbreak norm\_\allowbreak uebergang.py},
\texttt{k211\_\allowbreak symmetrie\_\allowbreak drehsinne.py} with their outputs, and the explorations in \texttt{notes/\allowbreak research/\allowbreak v4\_10\_0\_r03/};
pre-registration and result H94 in \texttt{notes/\allowbreak research/\allowbreak predictions/}.  \(\square\)

\paragraph{Consequence.}
Under G2 the big bang separates bosons from fermions.  Fermions are mirrored by CPT exactly and unitarily.  Bosons with matter weights are
Bogoliubov-transformed for \(k\Delta\ll1\) and CPT-mirrored for \(k\Delta\gg1\), and the transition is non-unitary in between, which the decision
\(\Delta\lesssim\ell_P\) moves to the Planck scale.  The two senses of rotation are exactly conjugate.  None of the tachyon-sector mechanisms yields a
vacuum term after the bang; \(\Lambda\) remains a free \(aV_{\min}\).  OP5 is not closed.

\subsection*{v4.10.0-r03 status in the three frames}
\label{subsec:v4100-r03-frames}

\begin{itemize}[leftmargin=2em]
\item r03.1--r03.4 concern the M--T transition under G2, with the weights \(a\), \(b\) of frame M: frame T before and frame M after the bang.  No frame-free
statement is made.
\item Photon frame P: with \(a=b\) there is no admixture, the fermion factor has modulus \(1\), and the pointwise identification of r03.6 is not obstructed.
The norm then stays \(+1\) for bosons, as in r85.5.
\item r03.5 states that frame T contributes no vacuum term in frame M after the bang under G2.  Under G1 both readings coexist, and the reduced theory is
the M theory with renormalized weights (r84).
\end{itemize}

\subsection*{v4.10.0-r03 register entries}
\label{subsec:v4100-r03-register}

\begin{center}
\small
\begin{tabular}{@{}>{\raggedright\arraybackslash}p{0.34\textwidth}>{\raggedright\arraybackslash}p{0.58\textwidth}@{}}
\toprule
statement & status after v4.10.0-r03\\
\midrule
fermions at the bang (G2) & no admixture, CPT exact, factor \(\pm i\beta_{\rm sym}\) (exploration)\\
bosonic norm at the bang & \(+1\to-1\) with one zero; non-unitary near \(k\sim1/\Delta\) (H94)\\
duration \(\Delta\) & \(\lesssim\ell_P\) (author decision)\\
senses of rotation & exactly conjugate, \(|C_+|=|C_-|\) (exploration)\\
\(\Lambda\) from I3, I4, ghost condensate under G2 & none after the bang; \(\Lambda=aV_{\min}\) free\\
pointwise identification of T and M & inconsistent for \(a\neq b\); boundary conditions withdrawn\\
\bottomrule
\end{tabular}
\end{center}

\subsection*{v4.10.0-r03 ledger and guard update}
\label{subsec:v4100-r03-ledger-guard}

\[
R_{\rm chk}^{v4100\text{-}r03}=(1,1,0)
\]
for (fermions and the bosonic norm at the bang settled; symmetry exact; \(\Lambda\) under G2 without a mechanism).

The layer spends no guard:
\[
R_{\rm glob}^{(v4100\text{-}r03,+)}=(0,1,1,1,1),
\qquad
\Gamma_{\rm guard}^{v4100\text{-}r03}=(1,1,1,1,1),
\qquad
\beta_4=1/256\ \text{unspent}.
\]
It does not imply exact finite \(C5\) zero closure, exact all-readout terminality,
\(C6/T2\), a global photon-clock/RH theorem, or global all-readout closure.

% ============================================================================
\clearpage
\section*{v4.10.0-r04 Source-Only: The Vacuum Term with a Tachyon Branch, and the Bosons Created at the Big Bang}
\addcontentsline{toc}{section}{v4.10.0-r04 Source-Only: The Vacuum Term with a Tachyon Branch, and the Bosons Created at the Big Bang}
\label{sec:v4100-r04-tachyon-branch}
% ============================================================================

\paragraph{Purpose.}
Layer v4.10.0-r03 found that under the second reading no tachyon-sector mechanism gives a vacuum term after the big bang
(\S\ref{sec:v4100-r03-bang-g2}).  This layer records the check under the domain G1, a temporary author decision to admit the tachyon branch
as degrees of freedom of its own, two analyses of the deferred idea I3 under that decision (fixed weights and a dynamic branch gas, H95),
the withdrawal of the decision, an estimate of the bosons created at the bang, and the state of the comparison data of the open predictions.
The notes are \texttt{notes/\allowbreak research/\allowbreak v4\_10\_0\_r04/}.  Earlier layers are not rewritten.

\noindent\textbf{Source-only status.}
This v4.10.0-r04 layer is source-only.  It does not build a PDF, does not assign a new DOI,
and does not authorize a public release.  The v4.10.0 Version DOI
\href{https://doi.org/10.5281/zenodo.23991357}{10.5281/zenodo.23991357} and the Concept DOI
\href{https://doi.org/10.5281/zenodo.19210728}{10.5281/zenodo.19210728} are unchanged by it.

\subsection*{v4.10.0-r04 results}
\label{subsec:v4100-r04-results}

\paragraph{Theoremlet v4.10.0-r04.1 (domain G1 gives no tension either).}
Exploration (analytic).  The mixture of I3 (k182, k183) consists of a normal branch (\(\omega=m\cosh\eta\), \(k=m\sinh\eta\)) and a tachyonic
branch (\(\omega=m\sinh\eta\), \(k=m\cosh\eta\)); the tension comes from the tachyonic one.  The reduced theory of the second reading,
\((a-b)\tfrac12X-(a+b)V\), has \(\omega^2=k^2+\gamma m^2\) for every mode and no tachyonic branch.  For a massive scalar the averaged pressure
vanishes (dust).  Under G1, as under G2, the only vacuum term is a free multiple of \(V_{\min}\).  I3 needs the tachyonic branch as degrees of freedom
of its own, that is, a reversal of the decision of r84.

\paragraph{Decision v4.10.0-r04.2 (author, 6 November 2026: tachyon branch admitted, later withdrawn).}
First the tachyonic branch was admitted as degrees of freedom of its own, so that I3 becomes possible.  The known cost: quanta of negative energy, and
vacuum decay through the graviton with a frame cut-off \(\lesssim14\) MeV (r83, H76).  After r04.3 and r04.4 the author withdrew the decision: the
tachyon part is again only the second reading of the same field (r84).

\paragraph{Theoremlet v4.10.0-r04.3 (I3 with fixed weights is excluded).}
Exploration (analytic, with an estimate).  The M and T readings of the mixture have the same density and dilute alike.  With fixed weights the sum is
one fluid with constant \(w=-b/(3a)\), not dust plus \(\Lambda\).  To explain \(\Omega_\Lambda\approx0.69\) today one needs \(b/a\approx2\), and then the dark
component at recombination is about \(10^{-5}\) of the baryons instead of \(5.4\) (cosmological values from prior knowledge, not looked up).  In
addition, the mixture at common rapidity is not preserved under redshift: \(k_1/k_2=\tanh\eta\) would stay fixed while each \(k\) is fixed by \(\eta\).

\paragraph{Theoremlet v4.10.0-r04.4 (H95: a dynamic branch gas cannot do it either).}
Pre-registered (commit \texttt{9eb4e40}, not blind; the refutation was expected analytically).  Model: flat FRW with radiation and baryons, and a
freely streaming gas with equal numbers in both branches.  The normal branch has weight \(1\); the tachyonic branch has weight \(r\), negative
energy, and threshold at \(a_{\rm th}=k_0/m\).  Grid \(60\times60\times40\) in \((r,a_{\rm th},\kappa_1/a_{\rm th})\): of \(144\,000\) points, \(12\,780\) keep
the total density positive for \(a\in[10^{-8},1]\).  None has both a dark-to-baryon ratio in \([4.32,6.48]\) at \(z=1090\) and \(w_{\rm eff}\in[-0.8,-0.6]\)
today.  \textbf{Refuted.}
\emph{Noticed afterwards.}  \(w_0\) can be made arbitrarily negative (down to \(-151\)), but only with the threshold practically today
(\(a_{\rm th}=1.000001\)), and then the ratio at recombination is \(266\).

\paragraph{Theoremlet v4.10.0-r04.5 (bosons created at the bang; estimate).}
Exploration.  With \(\Delta\lesssim\ell_P\) (r03.3), bosons with matter weights receive \(n_0=0.404\) per mode up to \(k_{\max}\sim1/\Delta\gtrsim M_P\).  This
gives \(\rho\approx g\,n_0\,k_{\max}^4/(8\pi^2)=g\cdot0.0051\,k_{\max}^4\), equivalent to a temperature of \(\approx0.35\,k_{\max}\).  The radiation after the bang is
Planck-hot, non-thermal, and initially without fermions.  The number depends on the non-unitary range \(k\sim1/\Delta\) of r03.2 and is an order of
magnitude only.

\paragraph{Remark v4.10.0-r04.6 (state of the comparison data).}
For H15 (\(m_W\)) no world average published after 2026-10-21 was found, and for H42/H46 (\(\theta_{23}\)) no global analysis after 2026-10-24; all three
remain open.  \emph{Disclosure:} The PDG 2026 edition (May 2026) already gave \(m_W=80.3625\pm0.0077\) GeV before the pre-registration of H15, which
listed only PDG 2025 as known.  The value is \(0.13\sigma\) from the prediction and is not a test.

\paragraph{Proof and measurement sources.}
Script \texttt{k212\_\allowbreak zweiggas\_\allowbreak frw.py} with output, and the explorations and the proposal in
\texttt{notes/\allowbreak research/\allowbreak v4\_10\_0\_r04/}; pre-registration and result H95 in \texttt{notes/\allowbreak research/\allowbreak predictions/}.  \(\square\)

\paragraph{Consequence.}
The deferred idea I3 (\(\Lambda\) from the dark-matter calculation) has failed in both forms examined, with fixed weights and with a dynamic branch gas,
even with the tachyon branch admitted.  The author returned to the second reading alone, so \(\Lambda\) is an input.  The bosons created at the bang
form a Planck-hot, non-thermal start without fermions.  An observable trace would have to survive thermalization, for example a gravitational-wave
background, which first needs a rule for how tensor modes see the transition.  OP5 is not closed.

\subsection*{v4.10.0-r04 status in the three frames}
\label{subsec:v4100-r04-frames}

\begin{itemize}[leftmargin=2em]
\item r04.1--r04.4 concern how the gravity of frame M sees frame T.  Under the second reading T is a re-weighting of the same field (no own branch).
With an own branch, T would have quanta of negative energy; that was examined and withdrawn.
\item r04.5 holds in frame M after the bang.  With photon weights (frame P, \(a=b\)) no bosons are created.
\item No frame-free statement is made.
\end{itemize}

\subsection*{v4.10.0-r04 register entries}
\label{subsec:v4100-r04-register}

\begin{center}
\small
\begin{tabular}{@{}>{\raggedright\arraybackslash}p{0.34\textwidth}>{\raggedright\arraybackslash}p{0.58\textwidth}@{}}
\toprule
statement & status after v4.10.0-r04\\
\midrule
\(\Lambda\) under G1 (second reading) & no tension; free multiple of \(V_{\min}\)\\
tachyon branch as own degrees of freedom & admitted, then withdrawn (author)\\
I3 with fixed weights & excluded (constant \(w\), CMB)\\
I3 with dynamic branch gas & refuted (H95)\\
bosons created at the bang & \(\rho\approx g\cdot0.0051\,k_{\max}^4\), non-thermal, no fermions (estimate)\\
H15, H42, H46 & open; PDG 2026 \(m_W\) disclosed as known before H15\\
\bottomrule
\end{tabular}
\end{center}

\subsection*{v4.10.0-r04 ledger and guard update}
\label{subsec:v4100-r04-ledger-guard}

\[
R_{\rm chk}^{v4100\text{-}r04}=(1,1,0)
\]
for (I3 excluded in both forms; decision withdrawn; \(\Lambda\) remains an input).

The layer spends no guard:
\[
R_{\rm glob}^{(v4100\text{-}r04,+)}=(0,1,1,1,1),
\qquad
\Gamma_{\rm guard}^{v4100\text{-}r04}=(1,1,1,1,1),
\qquad
\beta_4=1/256\ \text{unspent}.
\]
It does not imply exact finite \(C5\) zero closure, exact all-readout terminality,
\(C6/T2\), a global photon-clock/RH theorem, or global all-readout closure.

% ============================================================================
\clearpage
\section*{v4.10.0-r05 Source-Only: Gravitons at the Big Bang, a Review from the Core, and the Readout of \texorpdfstring{\(\Lambda\)}{Lambda}}
\addcontentsline{toc}{section}{v4.10.0-r05 Source-Only: Gravitons at the Big Bang, a Review from the Core, and the Readout of Lambda}
\label{sec:v4100-r05-core-review}
% ============================================================================

\paragraph{Purpose.}
This layer has three parts.  First, an estimate of gravitons created at the big bang and the author decision taken on it, with a correction to its
rationale.  Second, a review of the work since r81 from the core of the companion, prompted by the author's observation that the results keep
landing ``a little off'' while much structure appears.  The review leads to a separation of OP5 into three parts.  Third, the readout of \(\Lambda\)
at the present epoch, with one prediction without free parameter (H96) and one refuted dynamic bridge (H97).
The notes are \texttt{notes/\allowbreak research/\allowbreak v4\_10\_0\_r05/}.  Earlier layers are not rewritten.

\noindent\textbf{Source-only status.}
This v4.10.0-r05 layer is source-only.  It does not build a PDF, does not assign a new DOI,
and does not authorize a public release.  The v4.10.0 Version DOI
\href{https://doi.org/10.5281/zenodo.23991357}{10.5281/zenodo.23991357} and the Concept DOI
\href{https://doi.org/10.5281/zenodo.19210728}{10.5281/zenodo.19210728} are unchanged by it.

\subsection*{v4.10.0-r05 results}
\label{subsec:v4100-r05-results}

\paragraph{Theoremlet v4.10.0-r05.1 (gravitons at the big bang; estimate and correction).}
Exploration, not blind (the bound on \(N_{\rm eff}\) is known; cosmological values from prior knowledge).
\begin{enumerate}[leftmargin=2em]
\item[(a)] With matter weights, gravitons receive \(n_0=0.404\) per mode up to \(k_{\max}\sim1/\Delta\) (r03.3).  With the bosons of the Standard Model
thermalizing and the gravitons decoupled, this gives \(\Delta N_{\rm eff}\approx3.6\); the bound \(\approx0.3\) would need \(\gtrsim300\) bosonic degrees of freedom.
\item[(b)] Author decision (6 November 2026): gravity carries photon weights (\(a=b\)) at the transition.
\item[(c)] Correction to the rationale of (b): photon weights suppress creation only for \(k\Delta\ll1\).  At the smooth transition the produced energy
is of the same order (\(\int x^3n_{\rm mix}\,dx=0.066\) against \(0.091\,\Delta^{-4}\), maximum at \(k\Delta\approx0.76\)), so \(\Delta N_{\rm eff}\approx2.6\).
\item[(d)] Counting only modes below \(x_c=k_c\Delta\) keeps \(\Delta N_{\rm eff}\lesssim0.3\) for \(x_c\lesssim0.13\) (or \(\lesssim0.05\) if only the Higgs carries matter
weights).  This would be another chosen quantity.  \(\Delta N_{\rm eff}\) at the bang is therefore recorded as an open problem, not closed by a cut-off.
\end{enumerate}

\paragraph{Remark v4.10.0-r05.2 (review: why ``a little off'').}
A count of the tests since v4.8.0 (wiki table, after r04) gives about \(20\) confirmed, \(19\) refuted, \(9\) invalid, \(4\) open and \(28\) mixed results,
\(29\) marked not blind.  The confirmed tests without that mark all check inner mathematics (the OP5 comparison model, one construction).  Against
external measurements there is no confirmed blind prediction yet.
\begin{enumerate}[leftmargin=2em]
\item[(a)] Much of the structure that ``works'' is known, consistent physics taken up by the companion.  Examples: Dixon and Furey for one generation,
triality, Einstein--Cartan with the torsion \(3/16\), Boyle, Finn and Turok for the big bang, ghost condensation, WKB.
\item[(b)] What is new is the frame reading (M, T, P, \(\beta_{\rm sym}\), weights \(a\), \(b\), G2).  Its physical claims are identifications of single
dimensionless numbers at one point.  With many constants and combinations, agreement within a few percent is expected; r56 and r70.4 found
chance level.
\item[(c)] Since r81 about twelve author decisions were taken, most of them after a consequence did not fit.  Each kept a consequence in range; none
produced a confirmed new prediction.  The aim W7 (``vacuum term not free'') of the action programme is not met.
\end{enumerate}

\paragraph{Decision v4.10.0-r05.3 (author: OP5 separated into three parts).}
Reading the core shows two meanings of OP5.  In the architecture, OP5 is the reversal-stable return readout of the route whose source readout is
OP1; closed means \(\mathcal I(\mathrm{OP1})=\mathcal I(\mathrm{OP5}_{\rm sym})\) (Theorem~\ref{thm:op15-closure}), locally under the guards.  In physics,
\(\Lambda\) is a readout quantity, \(\Omega_\Lambda=1-\beta_{\rm sym}^2\) and \(\Omega_M=\beta_{\rm sym}^2\) as branching of the cascade
(Theorem~\ref{thm:dark-energy}), not a dynamical vacuum energy.  From late October the question became ``which dynamics produces \(\Lambda\)'' (OP5
series, then r81--v4.10.0-r04).  The negative results of that line answer a question the core does not ask.  Decision of 6 November 2026:
\begin{enumerate}[leftmargin=2em]
\item[(i)] \emph{OP5 (core):} the reversal closure of the source--return pair, locally closed, not claimed globally;
\item[(ii)] \emph{\(\Lambda\) readout:} whether \(\Omega_\Lambda:\Omega_M=(1-\beta^2):\beta^2\) is the cascade branching, and at which epoch;
\item[(iii)] \emph{action programme W1/W2:} a separate branch with the negative results of r81--v4.10.0-r04, frozen, without further repair decisions.
\end{enumerate}
A second decision fixes the epoch of (ii) as the present (way (a) of v4.8.0-r17).

\paragraph{Theoremlet v4.10.0-r05.4 (H96: \(\Omega_M\) today without free parameter; pre-registered, open).}
Pre-registered (commit \texttt{aeaba3d}, supplement on prior knowledge \texttt{e338b9b} before any comparison).  Prediction:
\[
\Omega_M=\beta_{\rm sym}^2=\frac{\gamma_L-1}{\gamma_L+1}=0.30185,\qquad \gamma_L=\gamma_{\rm BI}/\gamma_{\rm string}=1.864700.
\]
Comparison: the first flat-\(\Lambda\)CDM value of \(\Omega_m\) published after 2026-11-06 from DESI DR3 (BAO alone) or from Euclid; confirmed within
\(2\sigma\).  Known values (DESI DR2 \(0.2975\pm0.0086\), Planck \(0.3153\pm0.0073\) and others) are disclosed and do not count.  The prediction rests on the
U(1) counting of horizon states (r68) and on reading \(\gamma_L\), the ratio of a corrected to a superseded value of the same LQG parameter
(v4.8.0-r03 F11), as a Lorentz factor.

\paragraph{Theoremlet v4.10.0-r05.5 (the only dynamics for the ratio, and a bridge refuted).}
Search and one pre-registered test.
\begin{enumerate}[leftmargin=2em]
\item[(a)] The only place where the ratio follows from a dynamics is r31: the Markov cascade with return path has stationary flux fractions
\(\beta^2\), \(1-\beta^2\) for every rate.  A bridge to cosmic time is missing.  The earlier link of the core (scale factor as proper time, \(H\) as
activation rate of Riemann zeros) was withdrawn with the Riemann paper (v4.8.0-r04).
\item[(b)] Proposed bridge: the content of \(S^2\) is vacuum energy that returns to \(S^8\) at rate \(\Gamma=\varepsilon H\) and splits there.  Net, the vacuum
decays into matter at rate \(\beta^2\varepsilon H\); with \(\varepsilon=1\) this has no free parameter.
\item[(c)] H97 (commit \texttt{de36fcc}, not blind): this bridge with \(\varepsilon=1\) and \(\Omega_M=\beta^2\) against DESI DR2 BAO (13 values) gives
\(\chi^2_{\min}=105.0\) for \(12\) degrees of freedom, \(p=6\times10^{-17}\).  \textbf{Refuted.}  The control (flat \(\Lambda\)CDM, free \(\Omega_m\)) reproduces DESI
(\(\Omega_m=0.2970\), \(hr_d=101.56\) Mpc).  For comparison, flat \(\Lambda\)CDM at \(\Omega_M=\beta^2\) has \(\chi^2=10.9\), \(p=0.54\).
\emph{Exploration afterwards:} DR2 allows \(\varepsilon\lesssim0.4\) without preferring \(\varepsilon>0\).
\end{enumerate}

\paragraph{Proof and measurement sources.}
Scripts \texttt{k213\_\allowbreak gravitonen\_\allowbreak neff.py}, \texttt{k214\_\allowbreak erzeugung\_\allowbreak photongewichte.py},
\texttt{k215\_\allowbreak neff\_\allowbreak grenze.py}, \texttt{k216\_\allowbreak bruecke\_\allowbreak desi\_\allowbreak dr2.py} with outputs; the reviews and proposals in
\texttt{notes/\allowbreak research/\allowbreak v4\_10\_0\_r05/}; pre-registrations and results H96, H97 in \texttt{notes/\allowbreak research/\allowbreak predictions/}.  \(\square\)

\paragraph{Consequence.}
The companion is, at its core, a kinematic readout structure with sound algebra.  Its new physical claims are identifications.  They become
testable when they predict a function or several observables with one parameter, not one number at one point.  OP5 is separated accordingly.
The readout of \(\Lambda\) today gives a prediction without free parameter (H96, open).  The simplest dynamics for it is refuted (H97).
\(\Delta N_{\rm eff}\) at the bang is open.  OP5 (core) stays locally closed under the guards; no global closure is claimed.

\subsection*{v4.10.0-r05 status in the three frames}
\label{subsec:v4100-r05-frames}

\begin{itemize}[leftmargin=2em]
\item r05.1 concerns the M--T transition under G2; the photon-weight decision assigns gravity the weights of frame P.
\item r05.4 and r05.5: \(\beta_{\rm sym}^2\) is a matter-frame quantity (r22); the readout \(\Omega_M=\beta^2\) is made by a matter-frame observer.  In frame P it
vanishes (\(\Lambda_{\rm fundamental}=0\) of Theorem~\ref{thm:dark-energy}).  No frame-free statement is made.
\end{itemize}

\subsection*{v4.10.0-r05 register entries}
\label{subsec:v4100-r05-register}

\begin{center}
\small
\begin{tabular}{@{}>{\raggedright\arraybackslash}p{0.34\textwidth}>{\raggedright\arraybackslash}p{0.58\textwidth}@{}}
\toprule
statement & status after v4.10.0-r05\\
\midrule
gravitons at the bang & photon weights (author); \(\Delta N_{\rm eff}\) open (creation at \(k\sim1/\Delta\))\\
OP5 & separated: core closure (local) / \(\Lambda\) readout / action programme (frozen)\\
epoch of the \(\Lambda\) readout & present (author)\\
\(\Omega_M=\beta_{\rm sym}^2=0.30185\) & pre-registered prediction H96, open\\
cascade--Friedmann bridge, \(\varepsilon=1\) & refuted (H97)\\
blind predictions confirmed against data & none yet; open: H15, H42, H46, H96\\
\bottomrule
\end{tabular}
\end{center}

\subsection*{v4.10.0-r05 ledger and guard update}
\label{subsec:v4100-r05-ledger-guard}

\[
R_{\rm chk}^{v4100\text{-}r05}=(1,1,0)
\]
for (review done and OP5 separated; prediction H96 fixed; the dynamics for the readout absent).

The layer spends no guard:
\[
R_{\rm glob}^{(v4100\text{-}r05,+)}=(0,1,1,1,1),
\qquad
\Gamma_{\rm guard}^{v4100\text{-}r05}=(1,1,1,1,1),
\qquad
\beta_4=1/256\ \text{unspent}.
\]
It does not imply exact finite \(C5\) zero closure, exact all-readout terminality,
\(C6/T2\), a global photon-clock/RH theorem, or global all-readout closure.
% ============================================================================
\clearpage
\section*{v4.10.0-r06 Source-Only: Semiclassical Einstein Equations from Relative Entropy (Dorau--Much) and the Readouts of the Companion}
\addcontentsline{toc}{section}{v4.10.0-r06 Source-Only: Semiclassical Einstein Equations from Relative Entropy (Dorau--Much) and the Readouts of the Companion}
\label{sec:v4100-r06-relative-entropy}
% ============================================================================

\paragraph{Purpose.}
This layer takes up a recent result from the literature, pointed out by the author: P.~Dorau and A.~Much (Leipzig), \emph{From Quantum
Relative Entropy to the Semiclassical Einstein Equations}, Phys.\ Rev.\ Lett.\ \textbf{136}, 091602 (2026), DOI
\href{https://doi.org/10.1103/lmq8-nsty}{10.1103/lmq8-nsty}, arXiv:2510.24491.  It records what the paper shows and assumes, and where it meets the
readouts of the companion: the entropy factor of \S\ref{sec:bh} (v4.8.0-r03 F11a), the free \(\Lambda\) (v4.10.0-r03, r05), the sign of gravity
(v4.10.0-r01) and the frames.  The layer is an exploration of the literature; it contains no pre-registered test and no new prediction.
The note is \texttt{notes/\allowbreak research/\allowbreak v4\_10\_0\_r06/}.  Earlier layers are not rewritten.

\noindent\textbf{Source-only status.}
This v4.10.0-r06 layer is source-only.  It does not build a PDF, does not assign a new DOI,
and does not authorize a public release.  The v4.10.0 Version DOI
\href{https://doi.org/10.5281/zenodo.23991357}{10.5281/zenodo.23991357} and the Concept DOI
\href{https://doi.org/10.5281/zenodo.19210728}{10.5281/zenodo.19210728} are unchanged by it.

\subsection*{v4.10.0-r06 results}
\label{subsec:v4100-r06-results}

\paragraph{Theoremlet v4.10.0-r06.1 (the result of Dorau and Much, as stated by them).}
Literature, read in arXiv version v3; the formulas were not recomputed here.
\begin{enumerate}[leftmargin=2em]
\item[(a)] \emph{Setting.} Near a point the metric is approximated by \(\eta_{ab}\) (leading order).  The boost field \(\xi\) is null on two null
hypersurfaces meeting in a spacelike surface \(\mathcal S\): a local Rindler horizon, as in Jacobson's thermodynamic derivation (1995).  A real
Klein--Gordon field of mass \(m\) is quantized as a CCR algebra.  The reference state \(\omega_0\) is quasifree, boost invariant and Hadamard (in the
Rindler case the Minkowski vacuum); only coherent excitations \(\omega_\phi\) are treated.
\item[(b)] \emph{Relative entropy equals energy flux.} The modular group of the wedge acts as the boost, \(\Delta^{it}=\) boost by \(2\pi t\)
(Bisognano--Wichmann, Summers--Verch), and \(\omega_0\) is KMS at \(\beta=2\pi/\kappa\).  The Araki relative entropy then becomes
\[
S_{\rm rel}(\omega_0\Vert\omega_\phi)=-2\pi\int U\,\langle{:}T_{ab}{:}\rangle_{\omega_\phi}\xi^a\xi^b\,dU\,d\mathrm{vol}_{\mathcal S}.
\]
\item[(c)] \emph{Geometry.} With vanishing expansion, shear and twist at the horizon, the Raychaudhuri equation gives
\(\delta A=-\int U\,R_{ab}\xi^a\xi^b\,dU\,d\mathrm{vol}_{\mathcal S}\).
\item[(d)] \emph{Assumed, not derived:} \(S_{\rm rel}=\delta A/4\) (Bekenstein--Hawking, units \(G=\hbar=1\)).  This fixes the coupling \(\alpha=8\pi\).
\item[(e)] \emph{Result.} Holding for all null directions, \(\alpha\langle T_{ab}\rangle=R_{ab}+Ng_{ab}\); conservation and the contracted Bianchi identity
give \(N=-R/2+\Lambda\), hence \(G_{ab}+\Lambda g_{ab}=8\pi\langle{:}T_{ab}{:}\rangle\) with \(\Lambda\) an arbitrary integration constant.
\end{enumerate}
The authors call this ``arguments indicating'' the result: leading order, coherent states, scalar field.  The press reading ``gravity from vacuum
energy'' is not what is shown: the vacuum enters as the reference state of a relative entropy, not as a source of energy.

\paragraph{Theoremlet v4.10.0-r06.2 (the entropy factor becomes the gravitational coupling).}
Analytic, not blind (follows from r06.1 by one line).  Let a horizon counting give \(S=c\,A/4\) with \(c>0\).  Replacing (d) by
\(S_{\rm rel}=c\,\delta A/4\) and keeping the rest of r06.1 gives
\[
\alpha=\frac{8\pi}{c},\qquad G_{\rm eff}=\frac{G}{c}.
\]
For the two countings of \S\ref{sec:bh} (F11a in \S\ref{subsec:v480-r03-frame-parameter-provenance}) one has \(c=\gamma_*/\gamma\), where \(\gamma_*\) is the value
at which the counting reproduces \(S_{\rm BH}\): \(\gamma_{\rm string}=\ln2/(\pi\sqrt3)\) for the counting with \(j=\tfrac12\) only, \(\gamma_{\rm BI}=0.237533\) for
the counting of Domagala--Lewandowski and Meissner.  At one common value of \(\gamma\) the two countings differ by the factor
\(\gamma_{\rm BI}/\gamma_{\rm string}=\gamma_L=1.864700\) (F11a).  Under the chain of r06.1 this factor would appear as a ratio of gravitational couplings:
\[
\frac{G_{\rm eff}^{(j=1/2)}}{G_{\rm eff}^{(\rm DL/M)}}=\frac{c_{\rm DL/M}}{c_{j=1/2}}=\frac{\gamma_{\rm BI}}{\gamma_{\rm string}}=\gamma_L .
\]
\emph{Consequence.} If the T frame counts with \(j=\tfrac12\) and the M frame with all spins (r63, H47) at one common \(\gamma\), the two frames see different
Newton constants.  A single measured \(G\) is consistent only if each frame uses the \(\gamma\) at which its own counting reproduces \(S_{\rm BH}\); then
\(c=1\) in both, \(G\) is frame independent, and the frame difference sits in \(\gamma\).  This is the reading of r64 (the frame difference in \(\gamma\), not in
\(\hbar\) or \(G\)); Dorau--Much make it a statement about the coupling, not about the entropy alone.

\paragraph{Theoremlet v4.10.0-r06.3 (\(\Lambda\) as integration constant).}
Literature and comparison.  In r06.1(e) \(\Lambda\) is not fixed by the derivation; it enters through \(\nabla^a(\Lambda g_{ab})=0\).  This is the same situation
as in v4.10.0-r03 and r05: under G2 the action gives no \(\Lambda\) after the bang, \(\Lambda=a\,V_{\min}\) is free, and \(\Lambda\) is a readout
(Theorem~\ref{thm:dark-energy}).  The agreement supports the classification ``\(\Lambda\) is not part of the dynamics''.  It is not a success of the companion:
the standard derivation does not fix \(\Lambda\) either.  The readout \(\Omega_M=\beta_{\rm sym}^2\) (H96) is neither supported nor contradicted.

\paragraph{Remark v4.10.0-r06.4 (the sign of gravity from positivity; afterwards).}
In r06.1 the sign of \(\alpha\) is not free: \(S_{\rm rel}\ge0\) and \(c>0\) give \(\alpha>0\), i.e.\ attractive gravity.  In v4.10.0-r01 the sign of W1 was
corrected by an author decision after H90 had shown repulsion.  The corrected sign coincides with the one fixed by positivity.  This is noted afterwards
and is not a derivation of W1: Dorau--Much use a scalar field and Riemannian geometry, W1 is Einstein--Cartan with spinors and torsion
\(+(3/16)\kappa j^2\).  A transfer would need the relative entropy for Dirac fields, which the paper does not treat.

\paragraph{Remark v4.10.0-r06.5 (frames, Wick rotation and the modular group; open).}
The derivation uses that the modular group of the reference state is a real one-parameter group of boosts with KMS period \(2\pi\).  In the Euclidean
picture that period is the full rotation angle \(2\pi\) (the Unruh temperature via Wick rotation).  The companion has this bridge itself
(Theorem~\ref{thm:artqft}, \(t\mapsto i\tau_E\) on the biquaternion).  Under the tachyon map \(z\mapsto iz\) of Convention~F the rapidity becomes
imaginary and the boost a rotation.  Whether \(z\mapsto iz\) coincides with the Wick rotation of Theorem~\ref{thm:artqft} is open.  If it does, frame T is
the Euclidean side of the same KMS structure, not its loss.  The precise question is operator algebraic: is there a state on the wedge algebra whose
modular group is the tachyon image of the boost?  No answer is given here.

\paragraph{Proof and measurement sources.}
Exploration note \texttt{Erkundung\_\allowbreak Dorau\_\allowbreak Much\_\allowbreak relative\_\allowbreak Entropie.md} in
\texttt{notes/\allowbreak research/\allowbreak v4\_10\_0\_r06/}; the paper (arXiv:2510.24491v3).  r06.2 is one line of algebra on r06.1(d)--(e).  \(\square\)

\paragraph{Consequence.}
Dorau and Much give a rigorous quantum field theoretic form of Jacobson's argument.  They assume the area law and leave \(\Lambda\) free.  For the
companion this has three uses: the entropy factor of \S\ref{sec:bh} is the gravitational coupling (r06.2), so the frame reading of the horizon
countings must keep \(G\) frame independent; \(\Lambda\) as a readout is consistent with the standard derivation (r06.3); and the frames get a precise
operator algebraic question (r06.5).  Neither \(G\) nor \(\Lambda\) follows from it.

\subsection*{v4.10.0-r06 status in the three frames}
\label{subsec:v4100-r06-frames}

\begin{itemize}[leftmargin=2em]
\item M: the derivation of r06.1 holds as stated (real boosts, KMS at \(2\pi/\kappa\)).
\item T: open (r06.5); under \(z\mapsto iz\) the boost becomes a rotation, and whether the KMS structure survives as its Euclidean side is not decided.
\item P: degenerate fixed point; no separate statement.
\end{itemize}
r06.2 is frame dependent by construction (one counting per frame); r06.3 holds in every frame that has the derivation.  No frame-free statement is made.

\subsection*{v4.10.0-r06 register entries}
\label{subsec:v4100-r06-register}

\begin{center}
\small
\begin{tabular}{@{}>{\raggedright\arraybackslash}p{0.34\textwidth}>{\raggedright\arraybackslash}p{0.58\textwidth}@{}}
\toprule
statement & status after v4.10.0-r06\\
\midrule
Einstein equations from relative entropy & literature (Dorau--Much 2026); area law assumed, \(\Lambda\) free\\
entropy factor of \S\ref{sec:bh} & equals the gravitational coupling (r06.2); \(G\) frame independent only with matched \(\gamma\) per frame\\
\(\Lambda\) & integration constant in the standard derivation; readout in the companion (consistent)\\
sign of W1 (v4.10.0-r01) & coincides with the positivity sign (afterwards, scalar case only)\\
frame T and the modular group & open question (r06.5)\\
\bottomrule
\end{tabular}
\end{center}

\subsection*{v4.10.0-r06 ledger and guard update}
\label{subsec:v4100-r06-ledger-guard}

\[
R_{\rm chk}^{v4100\text{-}r06}=(1,1,0)
\]
for (literature recorded; entropy factor tied to the coupling; frame T under the modular group open).

The layer spends no guard:
\[
R_{\rm glob}^{(v4100\text{-}r06,+)}=(0,1,1,1,1),
\qquad
\Gamma_{\rm guard}^{v4100\text{-}r06}=(1,1,1,1,1),
\qquad
\beta_4=1/256\ \text{unspent}.
\]
It does not imply exact finite \(C5\) zero closure, exact all-readout terminality,
\(C6/T2\), a global photon-clock/RH theorem, or global all-readout closure.
% ============================================================================
\clearpage
\section*{v4.10.0-r07 Source-Only: The Relative-Entropy Chain in the Tachyon Frame, and a Correction to r06.5}
\addcontentsline{toc}{section}{v4.10.0-r07 Source-Only: The Relative-Entropy Chain in the Tachyon Frame, and a Correction to r06.5}
\label{sec:v4100-r07-tachyon-chain}
% ============================================================================

\paragraph{Purpose.}
The author asked whether the companion can map the derivation of Dorau and Much (\S\ref{sec:v4100-r06-relative-entropy}) onto its own frames and
so obtain the Einstein equations internally.  A pre-registered test was planned.  The search of prior work showed that the answer is fixed
analytically, so this layer is an exploration, not a test; no hypothesis number is assigned.  On the way, Remark r06.5 turned out to be wrong.  A second question of the author, why frame independence is reached and the prefactor of the area law is not, is answered in r07.4.
The notes are \texttt{notes/\allowbreak research/\allowbreak v4\_10\_0\_r07/}.  Earlier layers are not rewritten.

\noindent\textbf{Source-only status.}
This v4.10.0-r07 layer is source-only.  It does not build a PDF, does not assign a new DOI,
and does not authorize a public release.  The v4.10.0 Version DOI
\href{https://doi.org/10.5281/zenodo.23991357}{10.5281/zenodo.23991357} and the Concept DOI
\href{https://doi.org/10.5281/zenodo.19210728}{10.5281/zenodo.19210728} are unchanged by it.

\subsection*{v4.10.0-r07 results}
\label{subsec:v4100-r07-results}

\paragraph{Correction v4.10.0-r07.1 (to Remark r06.5).}
Remark v4.10.0-r06.5 states that under the tachyon map \(z\mapsto iz\) the rapidity becomes imaginary and the boost a rotation, and asks whether
\(z\mapsto iz\) is the Wick rotation of Theorem~\ref{thm:artqft}.  This is wrong.  With \(z=E+iB\) (v4.9.0-r02), \(z\mapsto iz\) is the
electromagnetic duality rotation \(E\mapsto-B\), \(B\mapsto E\), a map of field space, not of spacetime.  By Theoremlet v4.8.0-r20.1 it commutes with
every element of \(SO(3,\mathbb C)\), the proper Lorentz group acting on \(z\).  A boost stays a boost with real rapidity.  The question of r06.5 is
answered in the negative: \(z\mapsto iz\) is not the Wick rotation \(t\mapsto i\tau_E\); the two act on different objects.  The frame entries of r06 for
frame T are replaced by r07.2.

\paragraph{Theoremlet v4.10.0-r07.2 (the chain in frame T for the free Maxwell field).}
Exploration, not blind (fixed analytically before the check).  For the free Maxwell field, the chain of r06.1 gives in frame T the same Einstein
equation as in frame M: the same coupling, the same sign, the same inputs.
\begin{enumerate}[leftmargin=2em]
\item[(a)] The duality rotation \(U_\alpha\) is a symmetry of the free Maxwell field (conserved helicity).  It acts pointwise on \(F\), maps the algebra of
the Rindler wedge onto itself and leaves the vacuum invariant.
\item[(b)] It commutes with the boost (v4.8.0-r20.1), hence with the modular group of the vacuum (Bisognano--Wichmann).
\item[(c)] Relative entropy is invariant under such an automorphism: \(S_{\rm rel}(\omega_0\Vert\omega_\phi)=S_{\rm rel}(\omega_0\Vert\omega_{U\phi})\).
\item[(d)] The energy--momentum tensor is Hermitian in \(z\) (energy density \(\tfrac12|z|^2\), Poynting vector \((\bar z\times z)/2i\)) and therefore
duality invariant; the flux \(T_{ab}\xi^a\xi^b\) is unchanged.  The bilinear invariant \(s=z\cdot z\) changes sign: the Lagrangian density
\(\tfrac12(E^2-B^2)\) and \(E\cdot B\) are odd (v4.8.0-r06.3, v4.9.0-r40).
\end{enumerate}
\emph{Reading.} Gravity couples to the Hermitian tensor; the frames differ in the bilinear invariants.  For the Maxwell field, gravity is therefore
frame blind.  This turns the reading of v4.9.0-r64 (``gravity itself is frame independent'') into a statement for this case.  It is not new physics:
the duality invariance of the Maxwell tensor is textbook knowledge.

\paragraph{Proof and measurement sources.}
(a)--(c) are standard (duality symmetry of the free Maxwell field; invariance of the Araki relative entropy under automorphisms that fix the
reference state).  (d) and the commutation are checked numerically by \texttt{k217\_\allowbreak dualitaet\_\allowbreak gravitation.py} (2000 random fields and
boosts): commutation, invariance of the tensor and of the null flux, and the sign change of \(E^2-B^2\) and \(E\cdot B\) hold exactly in floating point;
a non-unitary scaling \((1+0.3i)z\) changes the tensor by \(O(1)\) (negative control).  \(\square\)

\paragraph{Theoremlet v4.10.0-r07.3 (bookkeeping: what the internal version gains).}
The chain is written in QAM bookkeeping for frames M and T, with the duality rotation as a declared map
(\texttt{qam\_\allowbreak dorau\_\allowbreak much\_\allowbreak companion.tex}); \texttt{qam.sty} raises no error.  The exported record gives for both frames the
Einstein equation with status \emph{conditional} and inputs \{local flatness, area law\}, the same inputs as in Dorau--Much.  The internal version
gains frame independence as a statement (r07.2).  It does not reduce the inputs: the area law, and with it the normalization \(1/4G\), remains an
input (r06.2).  For comparison, the same bookkeeping applied to the paper itself (\texttt{qam\_\allowbreak dorau\_\allowbreak much\_\allowbreak paper.tex}) reports a
promotion if the Einstein equation is declared proved, the area law as an input of the result, and, under the main-body typing of the
Bekenstein--Hawking law as a photon-frame statement (\S\ref{sec:bh}), a frame mixing without declared map.

\paragraph{Theoremlet v4.10.0-r07.4 (why frame independence is reached and the prefactor of the area law is not).}
Analytic, not blind (author question after r07.3).
\begin{enumerate}[leftmargin=2em]
\item[(a)] \emph{Frame independence is a statement about relations.}  The map between M and T is a symmetry of the system that fixes the reference state
and commutes with its modular group (r07.2), and every link of the chain is built from objects invariant under it.  Such an argument yields equalities
between the frames.  It says nothing about the value both frames share.
\item[(b)] \emph{The prefactor is an absolute number.}  The area law \(S=\delta A/4\) fixes the entropy per Planck area, i.e.\ the coupling (r06.2).  A
symmetry between frames cannot fix it; something must give one frame its absolute value.
\item[(c)] \emph{The countings used by the companion are calibrated to it.}  A horizon counting gives \(S=(\kappa/\gamma)\,A/\ell_P^2\) with a combinatorial
constant \(\kappa\); the value of the Barbero--Immirzi parameter at which it reproduces the area law \(S=A/(N\ell_P^2)\) is \(\gamma_*=N\kappa\).  Both literature
values used by the companion are of this kind: \(\gamma_{\rm string}=4\kappa_{\rm ABCK}\), \(\kappa_{\rm ABCK}=\ln2/(4\pi\sqrt3)\) (v4.8.0-r03 F11a), and
\(\gamma_{\rm BI}=4\kappa_{\rm DL/M}\).  Their ratio
\[
\gamma_L=\frac{\gamma_{\rm BI}}{\gamma_{\rm string}}=\frac{\kappa_{\rm DL/M}}{\kappa_{\rm ABCK}}=1.864700
\]
does not depend on \(N\): the prefactor cancels (checked for \(N=1,4,8\)).  So do \(\beta_{\rm sym}=(\gamma_L-1)/(\gamma_L+1)\) and \(\Omega_M=\beta_{\rm sym}^2\)
(H96).  Every quantity the companion builds from the two countings is exactly the part in which the normalization of the area law drops out.
\item[(d)] \emph{Consequence.}  With these inputs the prefactor cannot be derived, in principle and not only in practice: it was put into each
\(\gamma_*\) and removed again by taking the ratio.  The companion inherits the Barbero--Immirzi ambiguity of loop quantum gravity, in which the area
law is matched, not derived.  A derivation would need one absolute value of \(\gamma\) (or of the entropy per state) from the structure of the companion,
not a ratio; no such statement exists in the corpus.
\item[(e)] \emph{A caution.}  The state disc at the Hopf point has eigenvalues \(3/4\) and \(1/4\) (v4.9.0-r17).  Identifying that \(1/4\) with the one of the
area law would be a single-number identification of the kind recorded in v4.10.0-r05.2; there is no mechanism linking the two, and it is not proposed.
\end{enumerate}

\paragraph{Limits.}
Only the free Maxwell field.  For massive fields the tachyon image of the action is tachyonic (H71); a tachyonic field has no stable Hadamard
ground state and the premises of the chain fail.  Under the second reading (v4.9.0-r84, H79) the tachyon part is a weight on the matter action; how it
enters the flux is not examined.  Interacting fields: duality is in general not a symmetry.

\paragraph{Consequence.}
The companion reproduces the Dorau--Much derivation in both frames for the Maxwell field and so obtains frame-blind gravity as a statement.  It adds
no derivation of the area law; \(G\) and \(\Lambda\) remain as in r06.  The reason is structural (r07.4): the frame map gives relations,
and the countings enter only through a ratio in which the prefactor cancels.  The open front is one absolute value of \(\gamma\) from the structure.

\subsection*{v4.10.0-r07 status in the three frames}
\label{subsec:v4100-r07-frames}

\begin{itemize}[leftmargin=2em]
\item M: the chain of r06.1 holds as stated.
\item T: for the free Maxwell field the chain holds unchanged (r07.2); for massive fields the tachyon image has no stable ground state (open under
the second reading).
\item P: the null sector \(s=0\) is fixed by \(z\mapsto iz\); no separate statement.
\end{itemize}
For the Maxwell field, M and T agree: the Einstein equation of the chain is frame free in that case.

\subsection*{v4.10.0-r07 register entries}
\label{subsec:v4100-r07-register}

\begin{center}
\small
\begin{tabular}{@{}>{\raggedright\arraybackslash}p{0.34\textwidth}>{\raggedright\arraybackslash}p{0.58\textwidth}@{}}
\toprule
statement & status after v4.10.0-r07\\
\midrule
r06.5 (boost becomes rotation under \(z\mapsto iz\)) & wrong; \(z\mapsto iz\) commutes with \(SO(3,\mathbb C)\) (correction r07.1)\\
Einstein equation in frame T (Maxwell) & same as in frame M (exploration, not blind)\\
gravity frame independent (r64) & a statement for the free Maxwell field; open for massive fields\\
inputs of the internal chain & same as Dorau--Much (local flatness, area law)\\
prefactor of the area law & not derivable from the two countings: it cancels in \(\gamma_L\) (r07.4)\\
\bottomrule
\end{tabular}
\end{center}

\subsection*{v4.10.0-r07 ledger and guard update}
\label{subsec:v4100-r07-ledger-guard}

\[
R_{\rm chk}^{v4100\text{-}r07}=(1,1,0)
\]
for (r06.5 corrected; chain in frame T for the Maxwell field; area law still an input, prefactor cancels in the ratio of countings).

The layer spends no guard:
\[
R_{\rm glob}^{(v4100\text{-}r07,+)}=(0,1,1,1,1),
\qquad
\Gamma_{\rm guard}^{v4100\text{-}r07}=(1,1,1,1,1),
\qquad
\beta_4=1/256\ \text{unspent}.
\]
It does not imply exact finite \(C5\) zero closure, exact all-readout terminality,
\(C6/T2\), a global photon-clock/RH theorem, or global all-readout closure.
% ============================================================================
\clearpage
\section*{v4.10.0-r08 Source-Only: The Assumptions of the Companion and the Dependence of the Open Predictions}
\addcontentsline{toc}{section}{v4.10.0-r08 Source-Only: The Assumptions of the Companion and the Dependence of the Open Predictions}
\label{sec:v4100-r08-assumptions}
% ============================================================================

\paragraph{Purpose.}
Dorau and Much (\S\ref{sec:v4100-r06-relative-entropy}) start from few assumptions, set openly and motivated independently, and derive one
strict consequence.  This layer applies that way of working to the companion.  First it settles a suspected conflict between the second reading
and the photon weights of gravity.  Then it lists the assumptions of the companion with their origin, and it traces in QAM bookkeeping what each of
the four open predictions rests on.  The layer is an inventory; it contains no test and no new prediction.  The notes are
\texttt{notes/\allowbreak research/\allowbreak v4\_10\_0\_r08/}.  Earlier layers are not rewritten.

\noindent\textbf{Source-only status.}
This v4.10.0-r08 layer is source-only.  It does not build a PDF, does not assign a new DOI,
and does not authorize a public release.  The v4.10.0 Version DOI
\href{https://doi.org/10.5281/zenodo.23991357}{10.5281/zenodo.23991357} and the Concept DOI
\href{https://doi.org/10.5281/zenodo.19210728}{10.5281/zenodo.19210728} are unchanged by it.

\subsection*{v4.10.0-r08 results}
\label{subsec:v4100-r08-results}

\paragraph{Theoremlet v4.10.0-r08.1 (the source of gravity under the weights).}
Exploration, analytic.  The suspicion was: under the second reading the action is \((a-b)S_M\) (H79), and with photon weights \(a=b\) for gravity
(v4.10.0-r05.1(b)) the gravitating source of matter would vanish.
\begin{enumerate}[leftmargin=2em]
\item[(a)] The sum \((a-b)\) arises only where both readings live at the same place (domain G1, explored in v4.10.0-r04, not adopted).  Under the
adopted domain G2 the matter action is \(a\,S_M\) after the bang and \(-b\,S_M\) on its image before; no sum occurs.
\item[(b)] The photon weights concern the gravitational action: \(-b_g S_{\rm grav}\) before, \(a_g S_{\rm grav}\) after the bang, \(a_g=b_g\).  After the bang
\(a_g\tilde G=a\,\kappa\,T_M\), so \(G_{\rm eff}=(a/a_g)\,G\), one common factor absorbed into the measured \(G\).  There is no conflict.
\item[(c)] \(G_{\rm eff}\) is universal only if all non-gravitational fields carry the same weight \(a\) after the bang.  Different overall factors for two
species would give them different active gravitational masses at equal inertial mass.  This assumption (U) is used tacitly in v4.10.0-r03 and is
stated nowhere; it is recorded here.  The factors of H82 act within one action and change parameters, not universality.
\item[(d)] The limit of r07 for massive fields concerns only the domain before the bang: under G2 only the matter reading lives after it, so the chain
of r06.1 applies today in frame M to massive and massless fields alike.
\end{enumerate}

\paragraph{Remark v4.10.0-r08.2 (the list of assumptions).}
The assumptions are listed in \texttt{Annahmenliste\_\allowbreak Companion.md} with origin and one of six categories: \emph{V} set before the consequences
were known, with a reason of its own; \emph{R} set afterwards, a selection or repair after a consequence was known or did not fit; \emph{K} correction of an
error; \emph{A} working assumption; \emph{L} input from the literature; \emph{S} tacit.  An R assumption is not a flaw, but a prediction resting on it is less
blind by that selection.  Of the action programme W1/W2 (frozen since v4.10.0-r05), seven entries are R or A within two days; this is the repair series
recorded in r05.2.  None of them enters the open predictions.

\paragraph{Theoremlet v4.10.0-r08.3 (what the open predictions rest on).}
Bookkeeping in QAM (\texttt{qam\_\allowbreak annahmen\_\allowbreak vorhersagen.tex}; \texttt{qam.sty} raises no error).  All four open predictions have status
\emph{reading}: their weakest premises are readings of the core, not derivations.  Their inputs according to the exported record:
\begin{center}
\small
\begin{tabular}{@{}>{\raggedright\arraybackslash}p{0.28\textwidth}>{\raggedright\arraybackslash}p{0.64\textwidth}@{}}
\toprule
prediction & inputs\\
\midrule
H96, \(\Omega_M=\beta_{\rm sym}^2\) & \(\gamma_{\rm BI}\), \(\gamma_{\rm string}\), counting rule, epoch ``today''\\
H15, \(m_W\) & \(\gamma_{\rm BI}\), \(\gamma_{\rm string}\), counting rule, \(m_Z\), reading \(\cot\theta_W=\gamma_L\)\\
H46, \(\sin^2\theta_{23}\in[0.49,0.50]\) & neutrino reading, mask, neutrino data, \(\varepsilon=\eta\), \(\gamma_{\rm BI}\), \(\gamma_{\rm string}\), counting rule\\
H42, \(\sin^2\theta_{23}\le0.54\) & neutrino reading, mask, neutrino data\\
\bottomrule
\end{tabular}
\end{center}
\begin{enumerate}[leftmargin=2em]
\item[(a)] \emph{H15, H46 and H96 rest on the same core}: the two literature values, the assignment of counting rules to the frames (M all spins, T
\(j=\tfrac12\), U(1) rather than SU(2); v4.9.0-r63, r68) and, through the frame map, the reading of \(\gamma_L\) as Lorentz factor.  If the counting rule
fails, all three fail together.  With the SU(2) counting preferred in the literature (\(\gamma_L\approx2.15\)) one would get \(m_W=82.68\) GeV instead of
\(80.361\) GeV and \(\Omega_M=0.365\) instead of \(0.302\).  The U(1) rule was selected by the lepton window (r63) and the \(m_W\) relation was noticed
afterwards, so these two agreements are not independent of the selection.
\item[(b)] H42 and H46 share the neutrino reading (r55, r57); H46 adds \(\varepsilon=\eta\) and with it the \(\gamma_L\) core.
\item[(c)] Only H42 is independent of the \(\gamma_L\) core.  The four open predictions form two independent groups, not four.
\item[(d)] Forming \(\gamma_L\) without the declared map from T to M is reported by QAM as frame mixing.  \(\gamma_L\) is defined only through the reading it
is used to support.
\end{enumerate}

\paragraph{Proof and measurement sources.}
\texttt{Erkundung\_\allowbreak Quelle\_\allowbreak Gravitation\_\allowbreak Gewichte.md} (r08.1); \texttt{Annahmenliste\_\allowbreak Companion.md} (r08.2, r08.3);
\texttt{qam\_\allowbreak annahmen\_\allowbreak H96.tex} and \texttt{qam\_\allowbreak annahmen\_\allowbreak vorhersagen.tex} with exported records.  The SU(2) values in r08.3(a)
follow from \(m_W=m_Z\gamma_L/\sqrt{1+\gamma_L^2}\) and \(\Omega_M=(\gamma_L-1)/(\gamma_L+1)\).  \(\square\)

\paragraph{Consequence.}
The most loaded assumption of the companion is the assignment of counting rules to the frames.  Deriving it from the frames (OP1 since v4.9.0-r63)
would strengthen three open predictions at once; a refutation of one of them points first to that rule.  If H15, H46 and H96 are all confirmed, this
supports the core more than one hit and less than three independent ones.  Assumption U is now explicit.

\subsection*{v4.10.0-r08 status in the three frames}
\label{subsec:v4100-r08-frames}

\begin{itemize}[leftmargin=2em]
\item M: r08.1 concerns the matter reading after the bang; all four predictions are read out in frame M.
\item T: enters r08.3 through \(\gamma_{\rm string}\) (counting \(j=\tfrac12\)) and the map from T to M; under G2 the tachyon reading lives before the bang.
\item P: no separate statement.
\end{itemize}
No frame-free statement is made.

\subsection*{v4.10.0-r08 register entries}
\label{subsec:v4100-r08-register}

\begin{center}
\small
\begin{tabular}{@{}>{\raggedright\arraybackslash}p{0.34\textwidth}>{\raggedright\arraybackslash}p{0.58\textwidth}@{}}
\toprule
statement & status after v4.10.0-r08\\
\midrule
second reading vs.\ photon weights of gravity & no conflict under G2 (r08.1)\\
equal weight of all fields after the bang (U) & tacit assumption, now stated\\
open predictions H15, H42, H46, H96 & all status reading; H15, H46, H96 on one \(\gamma_L\) core, H42 independent\\
most loaded assumption & counting rule per frame (OP1)\\
\bottomrule
\end{tabular}
\end{center}

\subsection*{v4.10.0-r08 ledger and guard update}
\label{subsec:v4100-r08-ledger-guard}

\[
R_{\rm chk}^{v4100\text{-}r08}=(1,1,0)
\]
for (no conflict of weights; dependence of the predictions traced; the counting rule not derived).

The layer spends no guard:
\[
R_{\rm glob}^{(v4100\text{-}r08,+)}=(0,1,1,1,1),
\qquad
\Gamma_{\rm guard}^{v4100\text{-}r08}=(1,1,1,1,1),
\qquad
\beta_4=1/256\ \text{unspent}.
\]
It does not imply exact finite \(C5\) zero closure, exact all-readout terminality,
\(C6/T2\), a global photon-clock/RH theorem, or global all-readout closure.
% ============================================================================
\clearpage
\section*{v4.10.0-r09 Source-Only: The Counting Rule per Frame, Its Singular Position (H98) and a Principle}
\addcontentsline{toc}{section}{v4.10.0-r09 Source-Only: The Counting Rule per Frame, Its Singular Position (H98) and a Principle}
\label{sec:v4100-r09-counting-rule}
% ============================================================================

\paragraph{Purpose.}
Layer v4.10.0-r08.3 found that three of the four open predictions rest on one assumption: the assignment of horizon counting rules to the frames
(M counts all spins with weight \(2\), T counts \(j=\tfrac12\); v4.9.0-r63).  This layer asks how much the known agreements say about that assignment
(pre-registered test H98) and records a principle that selects it (exploration).  The notes are \texttt{notes/\allowbreak research/\allowbreak v4\_10\_0\_r09/}.
Earlier layers are not rewritten.

\noindent\textbf{Source-only status.}
This v4.10.0-r09 layer is source-only.  It does not build a PDF, does not assign a new DOI,
and does not authorize a public release.  The v4.10.0 Version DOI
\href{https://doi.org/10.5281/zenodo.23991357}{10.5281/zenodo.23991357} and the Concept DOI
\href{https://doi.org/10.5281/zenodo.19210728}{10.5281/zenodo.19210728} are unchanged by it.

\subsection*{v4.10.0-r09 results}
\label{subsec:v4100-r09-results}

\paragraph{Theoremlet v4.10.0-r09.1 (H98: the singular position of the pair).}
Pre-registered (note \texttt{Vorab\_\allowbreak Notiz\_\allowbreak H98\_\allowbreak Zaehlregel\_\allowbreak Sonderstellung.md}, commit \texttt{0cad2d8}); \emph{not blind with respect to
the data}: all three values were known, and so was the agreement of the chosen pair.  A rule is a spin range and a weight; \(\gamma\) solves
\(\sum_j g(j)\,e^{-2\pi\gamma a(j)}=1\) with \(a(j)=\sqrt{j(j+1)}\).  The fixed rule space \(S\) has the spin ranges all \(j\), half-integer \(j\), integer
\(j\ge1\), \(j=\tfrac12\) only, \(j=1\) only, and the weights \(2\) and \(2j+1\): nine distinct rules, \(36\) pairs with \(\gamma_M>\gamma_T\).  The three windows are the
lepton window \(\beta_{\rm sym}\in[0.545,0.570]\) (v4.9.0-r61), \(m_W\) within \(2\sigma\) of PDG 2025 (\(80.3692\pm0.0133\) GeV, via \(\cot\theta_W=\gamma_L\)), and
\(\Omega_M\) within \(2\sigma\) of DESI DR2 BAO (\(0.2975\pm0.0086\), via \(\Omega_M=\beta^2\)).
\begin{center}
\small
\begin{tabular}{@{}lc@{}}
\toprule
window & pairs inside\\
\midrule
lepton window & 2\\
\(m_W\) & 1\\
\(\Omega_M\) & 2\\
all three & 1: all spins with weight \(2\) for M, \(j=\tfrac12\) for T (DL/M with ABCK)\\
\bottomrule
\end{tabular}
\end{center}
\textbf{Confirmed} by the criterion.  Values of the pair: \(\gamma_L=1.864700\), \(\beta=0.54941\), \(m_W=80.3615\) GeV (\(0.6\sigma\)), \(\Omega_M=0.3018\) (\(0.5\sigma\)).  The
SU(2) pair of the literature (weight \(2j+1\) for M) gives \(\gamma_L=2.1515\) and misses all three windows.

\emph{Reading.}  Each single window admits other pairs; only \(m_W\) selects uniquely.  The \(36\) values of \(\gamma_L\) lie between \(1.006\) and \(3.51\); the \(m_W\)
window is about \(\pm0.003\) wide in \(\gamma_L\), so about \(0.08\) pairs would be expected inside it.  The weight of the result rests on the \(m_W\) relation,
which was itself found afterwards (H15 disclosure, look-elsewhere about \(0.17\)).  A chance triple agreement is of the order of a few per cent; this
is no significance in the strict sense.  H98 does not derive the assignment; it shows that among natural rules only this pair agrees with all three
known values.

\paragraph{Remark v4.10.0-r09.2 (a principle for the assignment; exploration, not blind).}
From the two exact lemmas of v4.9.0-r63 a principle P can be stated: a frame counts at a horizon exactly the states that
(P1) saturate the uncertainty bound, and (P2) are compatible with its energy condition, where M has all excitation levels and T relaxes to the floor and
reaches only the smallest non-zero excitation.  Energy and area are linked by the local first law of isolated horizons (Frodden, Ghosh and Perez,
2011): for an observer at fixed distance the energy is proportional to the area.  In the rule space \(S\), P1 excludes the weight \(2j+1\) except at
\(j=\tfrac12\) (where every state saturates, r63.2), P2 gives M all spins and T only \(j=\tfrac12\) (smallest area \(a(\tfrac12)=\sqrt3/2\)).  P selects exactly the
pair confirmed in r09.1.

The counting rule is thus not derived but reduced: instead of a selection by the data (category R in r08.2) it rests on two readings with a reason of
their own, uncertainty and energy, together with proved lemmas (r63.2, r63.3) and one literature input (the first law).  Both readings came from author
remarks after H47 and are not blind.  P1 is a reinterpretation: the counting of Meissner counts projections \(m\), not ``saturating states''
(v4.9.0-r68), and the SU(2) treatment violates P1.  The untested bridge is in P2: the particle energy of the tachyon at its floor (r63.3) is identified
with the horizon energy seen by an observer.  It could be examined through the boost energy at the Rindler horizon, the modular flow of
\S\ref{sec:v4100-r06-relative-entropy}.

\paragraph{Proof and measurement sources.}
\texttt{k218\_\allowbreak regelraum\_\allowbreak paare.py} with output (r09.1, sums to \(j_{\max}=4000\)); \texttt{Ergebnis\_\allowbreak H98\_\allowbreak Zaehlregel\_\allowbreak Sonderstellung.md};
\texttt{Erkundung\_\allowbreak Zaehlregel\_\allowbreak Prinzip.md} (r09.2, analytic on the rule space).  \(\square\)

\paragraph{Consequence.}
The assignment of counting rules, the most loaded assumption of the companion (r08.3), has a singular position among natural rules with respect to
three known values, and a principle with two readings selects it.  It is still an assumption.  A derivation needs the bridge of P2; a blind test needs a
consequence of P other than \(\gamma_L\).

\subsection*{v4.10.0-r09 status in the three frames}
\label{subsec:v4100-r09-frames}

\begin{itemize}[leftmargin=2em]
\item M: counts all spins with weight \(2\) (P1, P2: all levels).
\item T: counts \(j=\tfrac12\) only (P2: floor, smallest area).
\item P: carries no counting (v4.9.0-r63).
\end{itemize}
The rules differ between M and T by construction; no frame-free statement is made.

\subsection*{v4.10.0-r09 register entries}
\label{subsec:v4100-r09-register}

\begin{center}
\small
\begin{tabular}{@{}>{\raggedright\arraybackslash}p{0.34\textwidth}>{\raggedright\arraybackslash}p{0.58\textwidth}@{}}
\toprule
statement & status after v4.10.0-r09\\
\midrule
H98 (singular position of DL/M with ABCK) & confirmed, not blind with respect to the data\\
counting rule per frame (OP1) & reduced to two readings (P1, P2) and the local first law; not derived\\
bridge in P2 (floor energy vs.\ horizon energy) & open\\
\bottomrule
\end{tabular}
\end{center}

\subsection*{v4.10.0-r09 ledger and guard update}
\label{subsec:v4100-r09-ledger-guard}

\[
R_{\rm chk}^{v4100\text{-}r09}=(1,1,0)
\]
for (H98 confirmed; principle P recorded; the bridge of P2 not established).

The layer spends no guard:
\[
R_{\rm glob}^{(v4100\text{-}r09,+)}=(0,1,1,1,1),
\qquad
\Gamma_{\rm guard}^{v4100\text{-}r09}=(1,1,1,1,1),
\qquad
\beta_4=1/256\ \text{unspent}.
\]
It does not imply exact finite \(C5\) zero closure, exact all-readout terminality,
\(C6/T2\), a global photon-clock/RH theorem, or global all-readout closure.
% ============================================================================
\clearpage
\section*{v4.10.0-r10 Source-Only: The Bridge of the Counting Principle, and the Counting Rule as a Frame-Dependent Area Quantum}
\addcontentsline{toc}{section}{v4.10.0-r10 Source-Only: The Bridge of the Counting Principle, and the Counting Rule as a Frame-Dependent Area Quantum}
\label{sec:v4100-r10-area-quantum}
% ============================================================================

\paragraph{Purpose.}
Remark v4.10.0-r09.2 left one bridge open: the principle P lets the tachyon frame count only the smallest area excitation because it sits at its
floor (v4.9.0-r63.3), and links floor energy to horizon energy through the local first law.  This layer examines that bridge.  It does not hold in
that form; what remains is a sharper statement of the counting rule.  A parallel audit on GitHub reached the same verdict (r10.4).  Exploration, analytic, no criterion.  The notes are
\texttt{notes/\allowbreak research/\allowbreak v4\_10\_0\_r10/}.  Earlier layers are not rewritten.

\noindent\textbf{Source-only status.}
This v4.10.0-r10 layer is source-only.  It does not build a PDF, does not assign a new DOI,
and does not authorize a public release.  The v4.10.0 Version DOI
\href{https://doi.org/10.5281/zenodo.23991357}{10.5281/zenodo.23991357} and the Concept DOI
\href{https://doi.org/10.5281/zenodo.19210728}{10.5281/zenodo.19210728} are unchanged by it.

\subsection*{v4.10.0-r10 results}
\label{subsec:v4100-r10-results}

\paragraph{Theoremlet v4.10.0-r10.1 (the counting equation is an Unruh statistics, and the distance cancels).}
Analytic.  By the local first law (Frodden, Ghosh and Perez, 2013) an observer at proper distance \(\ell\) from the horizon assigns a puncture of spin
\(j\), area \(8\pi\gamma\,a(j)\) (\(\ell_P=1\)), the energy \(E_j=\gamma\,a(j)/\ell\), and has the Unruh temperature \(T_U=1/(2\pi\ell)\).  Hence
\[
E_j/T_U=2\pi\gamma\,a(j),
\]
independent of \(\ell\) (checked for \(\ell=0.5\) to \(1000\)).  The counting equation \(\sum_j g(j)\,e^{-2\pi\gamma a(j)}=1\) is this statistics (Ghosh and
Perez, 2011).  All static observers count alike: a higher or lower energy shifts \(E_j\) and \(T_U\) in the same ratio.

\paragraph{Theoremlet v4.10.0-r10.2 (the tachyon frame sees the same statistics).}
By v4.10.0-r07.2 the tachyon map commutes with the boost and fixes the vacuum; an observer in frame T sees the same vacuum with the same modular group,
hence the same KMS temperature.  The floor of r63.3 (particle energy \(E\to0\) at fixed invariant) is a property of a tachyon particle, not of the state an
observer sees at the horizon.  \textbf{Correction to Remark r09.2:} the bridge in P2 does not hold through the local first law and the Unruh temperature;
along that route M and T count alike.  ``T reaches only the smallest excitation'' would need a state of T that is not thermal with respect to the boost;
the vacuum is thermal.

\paragraph{Theoremlet v4.10.0-r10.3 (the counting rule as a frame-dependent area quantum).}
Analytic.  If both frames see the same statistics (r10.1, r10.2), the same geometry (r07.2) and both satisfy \(S=A/4\) (r06.2), the only remaining place for a
difference is the area quantum.  The assignment of counting rules to the frames (r08.3, r09) is then equivalent to a frame-dependent Barbero--Immirzi
parameter,
\[
\gamma_M=\gamma_{\rm BI},\qquad \gamma_T=\gamma_{\rm string},\qquad \gamma_M/\gamma_T=\gamma_L=1.864700:
\]
the smallest area quantum is smaller by the factor \(\gamma_L\) in frame T than in frame M, while the classical area and \(G\) are the same.  This is
v4.9.0-r64.1(b) (``the frame difference sits in \(\gamma\)'') in sharper form.  The duality rotation alone cannot change an area quantum, since it fixes
everything geometric; a reason would have to come from the boost part of the frame map (T as the floor of the mirror frame, r63.4), which only
reparametrizes the horizon.  No such reason is known.

\paragraph{Remark v4.10.0-r10.4 (relation to the parallel audit on GitHub).}
A source reconstruction of OP5 made with another language model (notes \texttt{op5\_\allowbreak spin\_\allowbreak selection\_\allowbreak audit\_\allowbreak 2026\_\allowbreak 11\_\allowbreak 06.md},
\texttt{op5\_\allowbreak common\_\allowbreak state\_\allowbreak access\_\allowbreak 2026\_\allowbreak 11\_\allowbreak 06.md}, \texttt{op5\_\allowbreak projector\_\allowbreak spin\_\allowbreak bridge\_\allowbreak 2026\_\allowbreak 11\_\allowbreak 06.md} in
\texttt{notes/\allowbreak research/}) reached \texttt{main} after layer r09.  It agrees with r10.2 and adds two points that bear on r09.2.
\begin{enumerate}[leftmargin=2em]
\item[(a)] Uncertainty saturation alone does not select \(j=\tfrac12\): \(|1,1\rangle\) and every \(|j,\pm j\rangle\) saturate as well.  ``Absolute minimum over \(j\)'' for T and
``extremal states at every \(j\)'' for M are two different rules; saturation is compatible with both and decides between neither.  P1 in r09.2 was used
only to exclude the weight \(2j+1\), which stands; the selection of \(j=\tfrac12\) for T rested on P2.
\item[(b)] The tachyon floor \(E=\mu/\sqrt{v^2-1}\to0\) does not contain \(j\); without a law linking \(\mu\) to the spin and a transition that changes \(j\), it gives no
relaxation to \(j=\tfrac12\).  This is the particle side of r10.2.
\item[(c)] A passive, invertible change of description cannot remove spin sectors; a T reading restricted to \(j=\tfrac12\) is either a restricted readout of the
same state (the rest stays in the balance) or a selection with its own weight.
\end{enumerate}
Together: P2 holds neither as an energy condition (r10.2) nor as a relaxation (b); P1 does not select the spin (a).  The singular position of r09.1 (H98)
is not affected; the assignment remains an assumption, now stated as a frame-dependent area quantum (r10.3).

\paragraph{Proof and measurement sources.}
\texttt{Erkundung\_\allowbreak Bruecke\_\allowbreak P2\_\allowbreak Unruh.md}; the cancellation of \(\ell\) in r10.1 is one line and was checked numerically.  \(\square\)

\paragraph{Consequence.}
Principle P keeps P1 (uncertainty saturation) as a reading; P2 as an energy condition is not supported along the route of the first law and should be
replaced by the plainer statement that each frame has its own area quantum.  That is an assumption, not a derivation.  The singular position found in
H98 is unaffected.  The three predictions on the \(\gamma_L\) core (H15, H46, H96) now visibly rest on one assumption: a frame-dependent area quantum with
ratio \(\gamma_L\).

\subsection*{v4.10.0-r10 status in the three frames}
\label{subsec:v4100-r10-frames}

\begin{itemize}[leftmargin=2em]
\item M and T: the same Unruh statistics and the same modular group (r10.1, r10.2); this part is frame free.
\item The area quantum differs: \(\gamma_M=\gamma_{\rm BI}\), \(\gamma_T=\gamma_{\rm string}\) (r10.3, assumption).
\item P: carries no counting.
\end{itemize}

\subsection*{v4.10.0-r10 register entries}
\label{subsec:v4100-r10-register}

\begin{center}
\small
\begin{tabular}{@{}>{\raggedright\arraybackslash}p{0.34\textwidth}>{\raggedright\arraybackslash}p{0.58\textwidth}@{}}
\toprule
statement & status after v4.10.0-r10\\
\midrule
bridge in P2 via first law and Unruh temperature & does not hold (correction to r09.2)\\
counting rule per frame & equivalent to a frame-dependent area quantum, \(\gamma_M/\gamma_T=\gamma_L\) (assumption)\\
reason for a frame-dependent area quantum & open\\
\bottomrule
\end{tabular}
\end{center}

\subsection*{v4.10.0-r10 ledger and guard update}
\label{subsec:v4100-r10-ledger-guard}

\[
R_{\rm chk}^{v4100\text{-}r10}=(1,1,0)
\]
for (bridge examined; counting rule restated; its reason open).

The layer spends no guard:
\[
R_{\rm glob}^{(v4100\text{-}r10,+)}=(0,1,1,1,1),
\qquad
\Gamma_{\rm guard}^{v4100\text{-}r10}=(1,1,1,1,1),
\qquad
\beta_4=1/256\ \text{unspent}.
\]
It does not imply exact finite \(C5\) zero closure, exact all-readout terminality,
\(C6/T2\), a global photon-clock/RH theorem, or global all-readout closure.

% ============================================================================
\clearpage
\section*{v4.10.0-r11 Source-Only: Integration of the OP5 Source Reconstruction Made with Another Language Model}
\addcontentsline{toc}{section}{v4.10.0-r11 Source-Only: Integration of the OP5 Source Reconstruction Made with Another Language Model}
\label{sec:v4100-r11-op5-reconstruction}
% ============================================================================

\paragraph{Purpose.}
On 6 and 7 November 2026 the author, working with another language model (Codex), traced the origin of OP5 back to the core and brought the result to
\texttt{main} (merge \texttt{4ae32c9}, branch \texttt{codex/\allowbreak op5-\allowbreak question-\allowbreak rebuild}): twelve notes
\texttt{notes/\allowbreak research/\allowbreak op5\_*\_\allowbreak 2026\_\allowbreak 11\_\allowbreak 06.md} and
\texttt{qsc\_\allowbreak back\_\allowbreak to\_\allowbreak basic\_\allowbreak idea\_\allowbreak 2026\_\allowbreak 11\_\allowbreak 06.md}, four QAM notes
\texttt{qam\_*\_\allowbreak 2026\_\allowbreak 11\_\allowbreak 0[67].md}, an optional QAM package \texttt{qam/\allowbreak tex/\allowbreak qam-physics.sty} with examples
and 34 tests, and wiki changes.  The work ran in parallel with layers v4.10.0-r06 to r10 and partly examined the same questions; layer r10.4 already cites three
of the notes.  This layer records the content, checks the load-bearing formulas with independent code, and places the findings against layers r05 to
r10, following the pattern of v4.9.0-r69.  It contains one author decision made through Codex, two precisions of layer r05 and one reading candidate.  No
test, no new hypothesis number.  The notes are \texttt{notes/\allowbreak research/\allowbreak v4\_10\_0\_r11/}.  Earlier layers are not rewritten.

\noindent\textbf{Source-only status.}
This v4.10.0-r11 layer is source-only.  It does not build a PDF, does not assign a new DOI,
and does not authorize a public release.  The v4.10.0 Version DOI
\href{https://doi.org/10.5281/zenodo.23991357}{10.5281/zenodo.23991357} and the Concept DOI
\href{https://doi.org/10.5281/zenodo.19210728}{10.5281/zenodo.19210728} are unchanged by it.

\subsection*{v4.10.0-r11 results}
\label{subsec:v4100-r11-results}

\paragraph{Decision v4.10.0-r11.1 (author, 6 November, through Codex: the historical OP5 derivation is suspended as a foundation).}
The historical chain was: spectral carrier \(\to\) trace formula \(\to\) metric action \(\to\) physical bare \(\Lambda\) term \(\to\) prescribed remainder \(\to\)
cancellation task (OP5, \texttt{sec:lambda\_cancel}).  The step to a well-defined metric action and its variation is missing, and the size of the bare term does
not follow (v4.8.0-r03 F3/F4, r04, r05).  The author decided to suspend this derivation as the active foundation and to rebuild the question from the
independently founded core.  New working question: \emph{which independently defined physical quantity can be obtained by an explicit, justified readout
of the algebraic core, and which checkable constraint follows, without presupposing the historical \(\Lambda\) cancellation?}  The decision is not a
refutation of the whole idea; earlier layers, negative results and pre-registered criteria stay.  H96 stays open, H97 refuted, W1/W2 frozen.  A
second precision of the author (7 November, Codex hand-over): the unity of M, T and P was always the hypothesis to be tested, not a presupposed outcome.

\paragraph{Theoremlet v4.10.0-r11.2 (what the reconstruction establishes; checked independently).}
Analytic; 31 checks with own code, all confirmed.
\begin{enumerate}[leftmargin=2em]
\item[(a)] \emph{Source--return package.}  Theorem~\ref{thm:op15-closure} holds as a conditional schema; in parts 21/22 no native route with evaluable
\(A_2,A_3,A_4\), class test and readout is instantiated.  The published mean \(m=(f+b)/2\) loses the branch datum \(d=(f-b)/2\): in the shear model on
\(\mathbb R^3\) with \(A_2(u,z,y)=(u,z+y,y)\), class \(u\) and readout \((u,z)\), the carriers \((1,2,\pm1)\) give the same \(m=(1,2)\) and forward values
\((1,3)\), \((1,1)\).  Class conservation with inverse transports does not fix the forward branch.
\item[(b)] \emph{Four denominators.}  The register ratio \(q=\lambda_-/\lambda_+=\tanh^2\eta\), the outflow fraction \(f_4/(f_4+f_2)\) from \(S^8\) in r31,
the contrast \((\gamma_M-\gamma_T)/(\gamma_M+\gamma_T)\) of r67 and \(\Omega_M\) (matter density over critical density) have different reference quantities.
In r31 the outflow fraction returns the input \(q\), while the occupation ratio \(p_4/p_2=qr\) depends on the return rate; r31 stabilises \(q\), it does
not derive it.  The smaller normalised eigenweight is \(q/(1+q)\), not \(q\).
\item[(c)] \emph{State selection.}  The family \(E_q=(1+q)^{-1/2}e_x\), \(B_q=(q/(1+q))^{1/2}e_y\) has \(\operatorname{tr}G=1\), \(E\cdot B=0\), the same
orientation and sector \(s_q=(1-q)/(1+q)>0\), and any ratio \(0<q<1\); in the \(\sigma_x\) analyser all have \(p_\pm=\tfrac12\).  The register algebra does
not select \(q\).  With positive energy weights the minus share is \(w_-q/(w_++w_-q)\); \(f_-=q\) needs \(w_-/w_+=1/(1-q)\).
\item[(d)] \emph{Spin selection.}  \(j(j+1)-m^2-|m|=(j-|m|)(j+|m|+1)\): saturation is \(|m|=j\) at every \(j\).  With weight 2, the spin set
\(\{\tfrac12\}\) gives \(\gamma=\ln2/(\pi\sqrt3)=0.127384\), the set \(\{\tfrac12,1\}\) gives \(0.201784\); saturation plus the counting equation fixes no
\(\gamma\) until the spin set is fixed.
\item[(e)] \emph{Common state and projector bridge.}  An invertible change of description cannot remove spin sectors; a restricted readout keeps the unread
rest in the balance, a selection has its own weight \(w\); ``not accessible'' does not imply ``zero energy''.  The chirality projector
\(C_{\rm ch}=\tfrac12(1+iL_{e_7})\) commutes with the tachyon map \(z\mapsto iz\) (r51, \texttt{Sektoren.lean}); a compatible injective bridge
\(B\) with \(C_{\rm ch}B=B\Pi\), \(U_DB=BU_{\rm spin}\) passes this on, \(\Pi U_{\rm spin}=U_{\rm spin}\Pi\).  No internal projector yields different M/T spin
masks without a declared mask-changing rule.
\item[(f)] \emph{Back to the basic idea.}  Under duality \(z\mapsto e^{i\alpha}z\) the invariant \(s=z\cdot z\) goes to \(e^{2i\alpha}s\), \(|s|\) is kept,
\(z\mapsto iz\) gives \(s\mapsto-s\); a configuration with \(s\neq0\) does not reach \(s=0\) under Lorentz and duality maps.  The compensator image
\(C(z)=\tfrac12(z+i\bar z)=(E+B)(1+i)/2\) has \(s'=2i|E'|^2\neq0\), it is not null.  ``One object'' has at least four meanings (same carrier, same physical
state, states linked by an active operation, different readouts of one state), which must be kept apart.
\end{enumerate}

\paragraph{Proof and measurement sources.}
\texttt{llm/\allowbreak k219\_\allowbreak nachrechnung\_\allowbreak codex\_\allowbreak op5.py}, output \texttt{k219\_\allowbreak out.txt} (groups A--H), protocol
\texttt{Integration\_\allowbreak Codex\_\allowbreak OP5\_\allowbreak Rekonstruktion.md}; the commutation in (e) is the existing Lean result.  \(\square\)

\paragraph{Remark v4.10.0-r11.3 (relation to layers r05 to r10).}
No contradiction.  In detail:
\begin{enumerate}[leftmargin=2em]
\item[(a)] \emph{Precision of r05.3(i).}  ``OP5 (core) locally closed'' is to be read as: closed as a conditional schema under the assumptions of
Theorem~\ref{thm:op15-closure}; a native instance of the route is open (r11.2(a)).  The guards are unaffected.
\item[(b)] \emph{Precision of r05.3(ii) and of the r05 note on sources.}  That note named r31 as the place where the ratio ``comes from a dynamics''.
Read narrowly, the dynamics gives stationarity; the branching number is an input (r11.2(b)).  The \(\Lambda\) readout is the additional identification
\(q=\Omega_M(\text{today})\), as pre-registered in H96.  This agrees with r08.2, where K1 (\(\Omega_M=\beta^2\) as cascade branching) is already listed as an
assumption.
\item[(c)] \emph{r08.3, r09.1.}  The \(\gamma_L\) core of H15, H46, H96 is the same chain; r11.2(b), (c) add the denominators and \(q\neq\lambda_-\).  The
freedom of r11.2(d) is exactly the rule space S of H98; the singular position found there is unaffected.
\item[(d)] \emph{r10.}  r11.2(d), (e) were taken up in r10.4.  The projector bridge confirms r10.3 from the internal side: as the duality rotation fixes
every geometric quantity, the tachyon map fixes the internal sectors; neither yields a frame-dependent mask or area quantum by itself.
\item[(e)] \emph{r03.2, r22, r24.}  r11.2(f) restates known results (the compensator image of a light wave is not a light wave; P is the degenerate fixed
point; R \(\neq\) P) and adds the shell argument.
\end{enumerate}

\paragraph{Remark v4.10.0-r11.4 (QAM as frame interface language).}
The package \texttt{qam-physics.sty} declares physics profiles (for example GR and QM), nested default contexts, implicit or explicit left roles, an exported
formula payload, and a conversion between profiles only with a declared map and an explicit target formula; 34 software cases and ten exports pass,
and the existing Lean checker accepts the records.  These are bookkeeping tests, not physics tests.  The author's precision: QAM is a frame interface
language; GR and its physics mapping hold in the frame assigned to them and are not carried to other frames automatically.  Relation: v4.10.0-r07.2 is
such a declared conversion (map \(z\mapsto iz\), commuting with \(SO(3,\mathbb C)\); target formula the same Einstein equation); the Dorau--Much example
does not identify the Rindler observer with M, T or P, as in r07.3.  In this language r10.3 says that the T profile differs from the M profile by its
declared area quantum, \(\gamma_T=\gamma_M/\gamma_L\), and that no map is known that produces it.

\paragraph{Remark v4.10.0-r11.5 (two observations from the r27 notes; one reading candidate).}
\begin{enumerate}[leftmargin=2em]
\item[(a)] With \(U_{\pm\alpha}\) on the Bloch disc, the direction mean is \(m=\cos2\alpha\,(x,z)\) and \((Q(f)+Q(b))/2=Q(m)+Q(d)\) for the second moment
\(Q(v)=vv^{\mathsf T}\).  At \(v=(1,0)\), \(\alpha=\pi/6\) the trace is \(1\) before and \(1/4\) after averaging; the missing \(3/4=\sin^2(\pi/3)\) is the squared
direction difference.  This is the factor \(3/4\) of r27 and of the reading ``rotate by \(\pi/6\), then compensate'' (r10, r11); consistent, not new.
\item[(b)] For a non-null field with flux \(\mathbf P=E\times B\), the observers at rapidity \(0\) and \(2\eta_{\rm ref}\) along \(\mathbf P\) see the same field
scale and opposite fluxes; the flux-free profile lies in the middle, and \(1/r=|s|^{-1}n=\cosh2\eta_{\rm ref}\) is the ordinary Lorentz factor between the
two (example \(E=2e_x\), \(B=e_y\): \(v=\tfrac12\) flux free, \(v=\tfrac45\) gives \(E'=2e_x\), \(B'=-e_y\)).  At \(r=1/\gamma_L\) the boost to the flux-free
profile has \(v=\beta_{\rm sym}\) (r17).  \emph{Reading candidate, a posteriori:} r10.3 asks for a reason from the boost part of the frame map; the pair of
flux-opposed observers at relative Lorentz factor \(\gamma_L\) is a concrete place where \(\gamma_L\) is an ordinary Lorentz factor.  Codex states expressly
that this does not identify the r27 step with an observer change, and nothing here identifies the pair with M and T.
\end{enumerate}

\paragraph{Consequence.}
The two lines of work agree.  OP5 in its historical form (cancellation of a bare \(\Lambda\)) is no longer an active task; what remains is the conditional
source--return schema (no native instance) and the additional identification \(q=\Omega_M\) of H96, which stays a pre-registered hypothesis without a
derivation of the bridge.  The open points are the same as after r10: the reason for the counting rule per frame (a frame-dependent area quantum), a
mechanism selecting \(q\), and a measurement map from register to density.  Codex proposes the meaning of ``one object'' as the next question before
any further spin or counting branch.

\subsection*{v4.10.0-r11 status in the three frames}
\label{subsec:v4100-r11-frames}

\begin{itemize}[leftmargin=2em]
\item M and T: the shell argument (r11.2(f)) links \(s>0\) and \(s<0\) of equal modulus by duality; the internal sectors are the same in both frames (r11.2(e));
frame free.
\item P: not on the orbit of \(s\neq0\); the compensator image is not P.  R \(\neq\) P (r24); \(\gamma_L\) compares M with R.
\item Conversions between frame profiles need a declared map and target formula (r11.4); the counting rule per frame remains an assumption (r10.3).
\end{itemize}

\subsection*{v4.10.0-r11 register entries}
\label{subsec:v4100-r11-register}

\begin{center}
\small
\begin{tabular}{@{}>{\raggedright\arraybackslash}p{0.34\textwidth}>{\raggedright\arraybackslash}p{0.58\textwidth}@{}}
\toprule
statement & status after v4.10.0-r11\\
\midrule
historical OP5 cancellation task & suspended as foundation (author, 6 November)\\
OP5 (core), r05.3(i) & conditional schema closed; native route open (precision)\\
\(\Lambda\) readout, r05.3(ii) & additional identification \(q=\Omega_M\) (H96, open); r31 stabilises, does not derive (precision)\\
selection of \(q\) by the register algebra & none (explicit family)\\
internal projector as M/T spin mask & no (commutes with the tachyon map)\\
\(\gamma_L\) as Lorentz factor of flux-opposed observers & algebraic fact; link to r10.3 a reading candidate\\
\bottomrule
\end{tabular}
\end{center}

\subsection*{v4.10.0-r11 ledger and guard update}
\label{subsec:v4100-r11-ledger-guard}

\[
R_{\rm chk}^{v4100\text{-}r11}=(1,1,0)
\]
for (reconstruction integrated; formulas checked; native route and selection of \(q\) open).

The layer spends no guard:
\[
R_{\rm glob}^{(v4100\text{-}r11,+)}=(0,1,1,1,1),
\qquad
\Gamma_{\rm guard}^{v4100\text{-}r11}=(1,1,1,1,1),
\qquad
\beta_4=1/256\ \text{unspent}.
\]
It does not imply exact finite \(C5\) zero closure, exact all-readout terminality,
\(C6/T2\), a global photon-clock/RH theorem, or global all-readout closure.
